\documentclass[10pt,a4paper]{article}
\usepackage{amsmath,amssymb,amsthm,graphicx}
\usepackage{amsmath,amssymb,amsthm,mathtools}
\usepackage[no-math]{fontspec}
\newenvironment{authorclaim}[1]{\par\medskip\noindent\textbf{#1.}\ \ignorespaces}{\par\medskip}
\allowdisplaybreaks
\title{Log-seed explicit PRG for read-once depth-two AC0 with parity via iterated restrictions}
\author{Dean Doron, Pooya Hatami, William M. Hoza}
\date{doi:10.4230/LIPIcs.CCC.2020.6; CC BY 3.0 Unported}
\begin{document}
\maketitle
\noindent\textit{Editorial scope.} Theorem1 main PRG mechanism counted once with complete Sections2–5 author mathematics/prose. Supporting tail/restriction lemmas and all seed/error/width/model extensions are not split. Standard prior components remain cited with complete local adaptations/sketches as supplied by authors. No reproof or author repair.\par\medskip
\noindent\textit{Allowed standard prerequisites (editorial).}\begin{itemize}
\item Small-bias/almost-independent constructions, Fourier norms, elementary symmetric tail bounds and prior PRG/ROBP components. Formal Theorem2[17], Lemma6[16,24], small-bias constructions[25,3], and cited Forbes–Kelley/CHRT/Meka–Reingold–Tal/Lee components remain standard cited prerequisites. All author-given local adaptations, proofs/sketches and every continuation are retained, including Lemmas27–29 as actually printed. Full33-entry original bibliography retained.
\end{itemize}\medskip
\section*{1.3 Log-Seed PRGs and Our Main Result}
The present work considers the question of whether the "iterated restrictions" approach can also yield a PRG with seed length $O(\log n)$ for some interesting class of tests. At first glance, this might seem doubtful, since after all we must pay for many pseudorandom restrictions. Nevertheless, we answer in the affirmative, proving the following theorem.
\begin{authorclaim}{Theorem1}
For all $n\in\mathbb N$ and $\epsilon>0$, there is an explicit $\epsilon$-PRG for read-once depth-2 $\mathrm{AC}^0[\oplus]$ formulas on $n$ input bits with seed length
\[
O(\log n)+\widetilde O(\log(1/\epsilon)).
\]
Specifically, the seed length of our PRG is $O(\log n+\log(1/\epsilon)\cdot(\log\log(1/\epsilon))^5)$. One can prove a lower bound of $\Omega(\log n+\log(1/\epsilon))$ on the seed length of any PRG for this model.\footnote{This lower bound holds already for fooling parity functions.}
\end{authorclaim}
\section*{1.4 Read-Once Depth-2 $\mathrm{AC}^0[\oplus]$ Formulas}
The class of functions that is fooled by our PRG (read-once depth-2 formulas over the basis $\{\wedge,\vee,\oplus\}$, with negations allowed at the inputs for free) is certainly of interest. It includes read-once CNFs and read-once $\mathbb F_2$-polynomials as special cases. The problems of fooling these classes have both received a lot of attention [12,16,4,23,24,22]. Previously, even for read-once CNFs, PRGs with seed length $O(\log n)$ were only known for constant error [8,12], whereas our PRG maintains seed length $O(\log n)$ with near-optimal error $\epsilon=\exp(-\widetilde\Omega(\log n))$. Meanwhile, for read-once $\mathbb F_2$-polynomials, no PRGs with seed length $O(\log n)$ were known at all prior to our work.
Gopalan et al. did give a PRG with near-optimal seed length $\widetilde O(\log(n/\epsilon))$ for read-once CNFs, and more generally for read-once depth-2 $\mathrm{AC}^0[\oplus]$ formulas with the property that the output gate is not $\oplus$ [16]. They used their PRG to construct a near-optimal hitting set for width-3 ROBPs [16]. A subsequent line of work provided near-optimal PRGs for all read-once depth-2 $\mathrm{AC}^0[\oplus]$ formulas [23,24,22].\footnote{The PRGs we are referring to were designed to fool read-once $\mathbb F_2$-polynomials, but in fact they fool all of read-once depth-2 $\mathrm{AC}^0[\oplus]$.}
Conversely, a read-once depth-2 $\mathrm{AC}^0[\oplus]$ formula can be simulated by a width-4 ROBP (after suitably permuting the variables). The problems of designing improved PRGs for width-4 ROBPs and for read-once $\mathrm{AC}^0[\oplus]$ formulas of any constant depth are two major frontiers in unconditional pseudorandomness [24,13]. The model we study in this paper is an interesting special case.

\section*{2 Subset-Wise Symmetric Polynomials}
In this section, we will formally define subset-wise symmetric polynomials and prove suitable tail bounds for them. This section can be read on its own, independent of the application to PRGs. We start by recalling known tail bounds for elementary symmetric polynomials.
\subsection*{2.1 Gopalan and Yehudayoff's Bounds for Symmetric Polynomials}
As a reminder, the $k$-th elementary symmetric polynomial is defined by
\[
S_k(y)=\sum_{\substack{I\subseteq[m]\\|I|=k}}\prod_{i\in I}y_i.
\]
We rely on the following tail bound by Gopalan and Yehudayoff [17]. As discussed in Section1.5.1, the bound says that if two $S_k$-s in a row are small, then all subsequent $S_k$-s are small.
\begin{authorclaim}{Theorem2 ([17])}
Let $y\in\mathbb R^m$, $\theta>0$, and $\ell\in\mathbb N$ satisfy $S_\ell^2(y)\le\theta^\ell/\ell!$ and $S_{\ell+1}^2(y)\le\theta^{\ell+1}/(\ell+1)!$. Then, for every $k\ge\ell$,
\[
|S_k(y)|\le\left(\frac{64e^2\theta\ell}{k}\right)^{k/2}.
\]
\end{authorclaim}
The exact statement of Theorem2 does not appear in Gopalan and Yehudayoff's work [17], but it follows readily from their analysis, and it was used previously by Meka et al. [24, Theorem5.2].\footnote{The careful reader will notice a slight discrepancy between the exact constants of Theorem2 on the one hand and the statements by Gopalan and Yehudayoff [17] and Meka et al. [24] on the other. This discrepancy reflects a minor mistake in the original paper by Gopalan and Yehudayoff [17] that we have here corrected.}
\subsection*{2.2 Our Tail Bounds for Subset-Wise Symmetric Polynomials}
Let $B=(B_1,\ldots,B_Q)$ be a partition of $[m]$, namely $[m]=B_1\sqcup\cdots\sqcup B_Q$. (The sets $B_1,\ldots,B_Q$ correspond to the "buckets" discussed in Section1.5.1.) Throughout this paper, let $\mathbb N$ denote the set of nonnegative integers, $\mathbb N=\{0,1,2,\ldots\}$. For a vector $\vec k=(\vec k[1],\ldots,\vec k[Q])\in\mathbb N^Q$ and $y\in\mathbb R^m$, we define the following polynomial:
\[
S_{\vec k,B}(y)=\sum_{\substack{I\subseteq[m]\\\forall q,\ |B_q\cap I|=\vec k[q]}}\prod_{i\in I}y_i.
\]
We name these polynomials as subset-wise symmetric polynomials, since for every $q\in[Q]$, $S_{\vec k}(y)$ when restricted to the $B_q$ variables is a degree $\vec k[q]$ symmetric polynomial.
Throughout this section we fix $B=(B_1,\ldots,B_Q)$ to be a partition of $[m]$. When the partition $B$ is clear from the context, we will simply write $S_{\vec k}$ instead of $S_{\vec k,B}$. To formulate our tails bounds for the subset-wise symmetric polynomials, we will need the following auxiliary polynomials:
\[
R_{\vec k}(y)\mathrel{\mathop:}=S_{\vec k}^2(y)\cdot\prod_{q=1}^Q\vec k[q]!.
\]
Given $c>1$, we will assign each vector $\vec k\in\mathbb N^Q$ a weight, defined as
\[
\|\vec k\|_{(c)}=\sum_{q=1}^Qc^q\vec k[q].
\]
(In our PRG application, $B_q$ will be the set of terms with approximately $c^q$ input variables, so $\|\vec k\|_{(c)}$ will be approximately the number of input variables participating in each monomial of $S_{\vec k}$, as outlined in Section1.5.1.) It is easy to verify that the above weight function is indeed a norm; however, we will not be using this observation.
The main result of this section is a tail-bound for subset-wise symmetric polynomials. In Lemma3, the parameter $A$ is analogous to the "cutoff point" $k_0$ discussed in Section1.5.1.
\begin{authorclaim}{Lemma3}
Suppose $c>1$ and $Q,A\in\mathbb N$ satisfy $A>\max\{(10^6c/(c-1))\cdot c^Q,2^{60}Q^2\}$. Let $Y$ be a random variable taking values in $\mathbb R^m$. Moreover, suppose for every $\vec k\in\mathbb N^Q$ with $\|\vec k\|_{(c)}\le A$,
\[
\mathbb E_Y[R_{\vec k}(Y)]\le2^{-\|\vec k\|_{(c)}/8}.
\]
Then, except with probability $2^{-A/2^{23}}$ over $y\sim Y$,
\[
\sum_{\substack{\vec k\in\mathbb N^Q\\\|\vec k\|_{(c)}>A}}|S_{\vec k}(y)|\le2^{-A/1024}.
\]
\end{authorclaim}
Lemma3 is similar in spirit to Theorem2: it says that if the "early" subset-wise symmetric polynomials are small (with high probability), then the "late" subset-wise symmetric polynomials are all small (with high probability).
\subsection*{2.3 Non-probabilistic Tail Bound}
Before moving to the proof of Lemma3 in the next subsection, here we first give a tail-bound in the case when the input $y$ satisfies some useful properties. We will later prove Lemma3, by showing that a random $Y$ satisfies these properties with high probability. Given a vector $\vec k\in\mathbb N^Q$, we define the restriction of $\vec k$ to a set $\mathcal Q\subseteq[Q]$ by
\[
\vec k|_{\mathcal Q}[q]=\begin{cases}\vec k[q]&q\in\mathcal Q,\\0&q\notin\mathcal Q.\end{cases}
\]
Our non-probabilistic tail bound goes as follows.
\begin{authorclaim}{Lemma4}
Suppose $c>1$ and $Q,A\in\mathbb N$ satisfy $A>\max\{(10^6c/(c-1))\cdot c^Q,2^{60}Q^2\}$. Let $y\in\mathbb R^m$ be a fixed vector. Suppose that for every $\vec k\in\mathbb N^Q$, with $A/10^5\le\|\vec k\|_{(c)}\le A$, and for every pair of disjoint sets $\mathcal Q_1,\mathcal Q_2\subseteq[Q]$ satisfying $\{q:\vec k[q]>1\}\subseteq\mathcal Q_1\cup\mathcal Q_2$, we have
\[
R_{\vec k|_{\mathcal Q_1}}(y)\cdot R_{\vec k|_{\mathcal Q_2}}(y)^4\le2^{-\|\vec k\|_{(c)}/32}.
\]
Then,
\[
\sum_{\substack{\vec k\in\mathbb N^Q\\\|\vec k\|_{(c)}>A}}|S_{\vec k}(y)|\le2^{-A/1024}.
\]
\end{authorclaim}
\begin{proof}
For a fixed $\ell\in\mathbb N$ and $q\in[Q]$, define
\[
S_{\ell,q}=\sum_{\substack{I\subseteq B_q\\|I|=\ell}}\prod_{i\in I}y_i,
\]
which is the $\ell$-th elementary symmetric polynomial applied to $(y_i)_{i\in B_q}$. Similarly, define $R_{\ell,q}=S_{\ell,q}^2(y)\cdot\ell!$.
Fix $\vec k$ with $\|\vec k\|_{(c)}>A$, let $\lambda=10^5\cdot\|\vec k\|_{(c)}/A$ and let $\vec k'\in\mathbb N^Q$ be such that $\vec k'[q]=\lceil\vec k[q]/\lambda\rceil$. Thus, $A/10^5\le\|\vec k'\|_{(c)}\le A/2$. Let $\mathcal Q:=\{q\in Q:\vec k[q]\ge1\}$, and for each $q\in\mathcal Q$, let $\theta_q>0$ be the smallest\footnote{It is possible that $\theta_q=0$ satisfies Equation(3). In this degenerate case, we must have $S_{\vec k[q],q}=0$. This implies $S_{\vec k}(y)=0$, hence Equation(4) trivially holds.} number satisfying
\begin{equation}
R_{\vec k'[q],q}\le\theta_q^{\vec k'[q]}
\quad\text{and}\quad
R_{\vec k'[q]+1,q}\le\theta_q^{\vec k'[q]+1}.\tag{3}
\end{equation}
By Theorem2,
\[
|S_{\vec k[q],q}|\le\left(\frac{64e^2\theta_q\vec k'[q]}{\vec k[q]}\right)^{\vec k[q]/2}.
\]
Subset-wise symmetric polynomials by design can be expressed as a product of elementary symmetric polynomials, hence
\begin{align*}
|S_{\vec k}(y)|&=\prod_{q=1}^Q|S_{\vec k[q],q}(y)|
\le\prod_{q\in\mathcal Q}\left(\frac{64e^2\theta_q\vec k'[q]}{\vec k[q]}\right)^{\vec k[q]/2}\\
&=\left(\prod_{q\in\mathcal Q}\theta_q^{\vec k[q]/\lambda}\right)^{\lambda/2}
\cdot\prod_{q\in\mathcal Q}\left(8e\sqrt{\vec k'[q]/\vec k[q]}\right)^{\vec k[q]}.
\end{align*}
By our choice of $\theta_q$,
\begin{align*}
\theta_q^{\vec k'[q]}
&=\max\{R_{\vec k'[q],q}(y),R_{\vec k'[q]+1,q}(y)^{\vec k'[q]/(\vec k'[q]+1)}\}\\
&\le\max\{R_{\vec k'[q],q}(y),R_{\vec k'[q]+1,q}(y),\sqrt{R_{\vec k'[q]+1,q}(y)}\}.
\end{align*}
Observe that $\vec k[q]/\lambda\in[\vec k'[q]-1,\vec k'[q]]$, and thus $\theta_q^{\vec k[q]/\lambda}$ is between $\theta_q^{\vec k'[q]-1}$ and $\theta_q^{\vec k'[q]}$. If $\vec k'[q]=1$, then $\theta_q^{k'_q-1}=1$, and otherwise $\theta_q^{\vec k'[q]-1}$ is between $\theta_q^{\vec k'[q]}$ and $\sqrt{\theta_q^{\vec k'[q]}}$. Therefore,
\begin{align*}
\theta_q^{\vec k[q]/\lambda}
&\le\max\{\theta_q^{\vec k'[q]},\sqrt{\theta_q^{\vec k'[q]}},1_{\vec k'[q]=1}\}\\
&\le\max\{R_{\vec k'[q],q}(y),R_{\vec k'[q],q}(y)^{1/4},R_{\vec k'[q]+1,q}(y),R_{\vec k'[q]+1,q}(y)^{1/4},1_{\vec k'[q]=1}\}.
\end{align*}
For every $q$, choose $\vec k''[q]\in\{\vec k'[q],\vec k'[q]+1\}$ such that
\[
\theta_q^{\vec k[q]/\lambda}\le\max\{R_{\vec k''[q],q}(y),R_{\vec k''[q],q}(y)^{1/4},1_{\vec k''[q]=1}\}.
\]
Note that $\|\vec k''\|_{(c)}\ge\|\vec k'\|_{(c)}$ and
\[
\|\vec k''\|_{(c)}\le\|\vec k'\|_{(c)}+\sum_{q=1}^Qc^q\le\|\vec k'\|_{(c)}+\frac A{10^6}<A.
\]
Therefore, there exist disjoint sets $\mathcal Q_1,\mathcal Q_2\subseteq[Q]$ such that $\{q:\vec k''[q]>1\}\subseteq\mathcal Q_1\cup\mathcal Q_2$, and that for every $q\in\mathcal Q$,
\[
\theta_q^{\vec k[q]/\lambda}\le
\begin{cases}R_{\vec k''[q],q}(y)&q\in\mathcal Q_1,\\R_{\vec k''[q],q}(y)^{1/4}&q\in\mathcal Q_2,\\1&\text{otherwise}.\end{cases}
\]
Multiplying over $q\in\mathcal Q$, we get
\begin{align*}
\prod_{q\in\mathcal Q}\theta_q^{\vec k[q]/\lambda}
&\le\prod_{q\in\mathcal Q_1}R_{\vec k''[q],q}(y)\cdot\prod_{q\in\mathcal Q_2}R_{\vec k''[q],q}(y)^{1/4}\\
&=\left(R_{\vec k''|_{\mathcal Q_1}}(y)^4\cdot R_{\vec k''|_{\mathcal Q_2}}(y)\right)^{1/4}
\le2^{-\|\vec k''\|_{(c)}/128}\le2^{-\|\vec k'\|_{(c)}/128}.
\end{align*}
As a result,
\begin{align}
|S_{\vec k}(y)|
&\le2^{-\frac{\|\vec k'\|_{(c)}}{128}\cdot\frac\lambda2}
\cdot\prod_{q\in\mathcal Q}\left(8e\sqrt{\vec k'[q]/\vec k[q]}\right)^{\vec k[q]}\notag\\
&\le2^{-\frac{\|\vec k'\|_{(c)}}{256}\lambda}\cdot
\prod_{q\in\mathcal Q}\left(8e\sqrt{2/10^5}\right)^{\vec k[q]}\cdot\left(\sqrt{10^5/2}\right)^\lambda\notag\\
&\le2^{-\frac{\|\vec k'\|_{(c)}}{256}\lambda}\cdot2^{8Q\lambda}\cdot4^{-\|\vec k\|_1}
\le2^{-\frac{\|\vec k'\|_{(c)}}{512}\lambda}\cdot4^{-\|\vec k\|_1}
\le2^{-\|\vec k\|_{(c)}/512}\cdot4^{-\|\vec k\|_1}.\tag{4}
\end{align}
To see the second inequality, observe that when $\vec k[q]>\lambda$, then
$\left(8e\sqrt{\vec k'[q]/\vec k[q]}\right)^{\vec k[q]}
\le(8e\sqrt{2/10^5})^{\vec k[q]}$, and otherwise
$\left(8e\sqrt{\vec k'[q]/\vec k[q]}\right)^{\vec k[q]}\le(8e)^\lambda$.
Summing up over all choices of $\vec k$ we get,
\begin{align*}
\sum_{\substack{\vec k\in\mathbb N^Q\\\|\vec k\|_{(c)}>A}}|S_{\vec k}(y)|
&\le\sum_{L=1}^m\sum_{\substack{\vec k\in\mathbb N^Q\\\|\vec k\|_{(c)}>A,\|\vec k\|_1=L}}2^{-\|\vec k\|_{(c)}/512}\cdot4^{-L}\\
&\le2^{-A/512}\cdot\sum_{L=1}^m4^{-L}\cdot|\{\vec k\in\mathbb N^q:\|\vec k\|_1=L\}|\\
&=2^{-A/512}\cdot\sum_{L=1}^m4^{-L}\cdot\binom{Q-1+L}{Q-1}\\
&\le2^{-A/512}\cdot\sum_{L=1}^m4^{-L}\cdot2^{Q-1+L}
\le2^{-A/512}\cdot2^{Q-1}\cdot\sum_{L=1}^m2^{-L}\le2^{-A/1024}.
\end{align*}
\end{proof}
\subsection*{2.4 Probabilistic Tail Bound: Proof of Lemma3}
\begin{proof}
Let $\vec k,\mathcal Q_1$, and $\mathcal Q_2$ be as in the statement of Lemma4. Using the Cauchy-Schwarz inequality and the concavity of $(\cdot)^{1/4}$, we get
\begin{align*}
\mathbb E[(R_{\vec k|_{\mathcal Q_1}}(Y))^{1/8}(R_{\vec k|_{\mathcal Q_2}}(Y))^{1/2}]
&\le\left(\mathbb E[(R_{\vec k|_{\mathcal Q_1}}(Y))^{1/4}]\cdot\mathbb E[R_{\vec k|_{\mathcal Q_2}}(Y)]\right)^{1/2}\\
&\le\left(\mathbb E[R_{\vec k|_{\mathcal Q_1}}(Y)]^{1/4}\cdot\mathbb E[R_{\vec k|_{\mathcal Q_2}}(Y)]\right)^{1/2}\\
&\le(2^{-\|\vec k|_{\mathcal Q_1}\|_{(c)}/32}\cdot2^{-\|\vec k|_{\mathcal Q_2}\|_{(c)}/8})^{1/2}\\
&\le2^{-\|\vec k|_{\mathcal Q_1\cup\mathcal Q_2}\|_{(c)}/64}\\
&\le2^{-(\|\vec k\|_{(c)}-(c/(c-1))c^Q)/64}\\
&\le2^{-(\|\vec k\|_{(c)}-A/20000)/64}\le2^{-\|\vec k\|_{(c)}/128}.
\end{align*}
Therefore, by Markov's inequality, except with probability at most $2^{-\|\vec k\|_{(c)}/256}\le2^{-A/2560000}$, we have
\[
R_{\vec k|_{\mathcal Q_1}}(Y)\cdot R_{\vec k|_{\mathcal Q_2}}(Y)^4\le2^{-\|\vec k\|_{(c)}/32}.
\]
The above analysis was done for a fixed choice of $\vec k,\mathcal Q_1$, and $\mathcal Q_2$. The number of choices for such $\vec k$ is $A^Q$ (which is subexponential in $A$), and the number of such $\mathcal Q_1,\mathcal Q_2$ is at most $3^Q$ (which is a polynomial in $A$), thus Lemma3 follows by a union bound. More precisely, one can check that since $A\ge2^{60}Q^2$, then $(3A)^Q\cdot2^{-A/2560000}\le2^{-A/2^{23}}$.
\end{proof}

\section*{3 Pseudorandomness Preliminaries}
Having completed our analysis of subset-wise symmetric polynomials, we now move on to setting the groundwork for our PRG construction and analysis.
\subsection*{3.1 Probability Basics}
Let $U_n$ denote the uniform distribution over $\{0,1\}^n$. We will simply write $U$ if $n$ is clear from context. For $f:\{0,1\}^n\to\mathbb R$, as a shorthand, we write $\mathbb E[f]$ to denote $\mathbb E[f(U)]$ and $\operatorname{Var}[f]$ to denote $\operatorname{Var}[f(U)]$. If $X$ is a distribution over $\{0,1\}^n$, we say that $X$ $\epsilon$-fools $f$, or $X$ fools $f$ with error $\epsilon$, if $|\mathbb E[f(X)]-\mathbb E[f]|\le\epsilon$. We say that $X$ $\epsilon$-fools a family $\mathcal F$ of functions, if it $\epsilon$-fools every $f\in\mathcal F$.
\subsection*{3.2 Small Bias}
A parity function is a function of the form $f(x)=\bigoplus_{i\in I}x_i$ for some set $I\subseteq[n]$. We say that a random variable $Y\in\{0,1\}^n$ is $\delta$-biased if it $\delta$-fools all parity functions. We say that $Y$ is $n'$-wise $\delta$-biased if it $\delta$-fools all parity functions on at most $n'$ bits, i.e., all parity functions with $|I|\le n'$. There are explicit constructions of $n'$-wise $\delta$-biased distributions that can be sampled with $O(\log(n'/\delta)+\log\log n)$ truly random bits [25,3].
Recall that for a function $f:\{0,1\}^n\to\mathbb R$ with Fourier expansion $f=\sum_{S\subseteq[n]}\widehat f(S)\cdot\chi_S$, the $L_1$ norm of $f$ is defined by $L_1(f)=\sum_{S\subseteq[n]}|\widehat f(S)|$. This norm is subadditive ($L_1(f+g)\le L_1(f)+L_1(g)$) and submultiplicative ($L_1(f\cdot g)\le L_1(f)\cdot L_1(g)$). Functions with bounded $L_1$ norm are fooled by small-bias distributions:
\begin{authorclaim}{Claim5}
If $f:\{0,1\}^n\to\mathbb R$ and $Y$ is $\delta$-biased, then $Y$ fools $f$ with error $2\delta\cdot L_1(f)$.
\end{authorclaim}
We will also rely on the following "XOR lemma" for small-bias distributions.
\begin{authorclaim}{Lemma6 ([16,24])}
Let $0<\delta<\epsilon\le1$. Let $f_1,\ldots,f_k:\{0,1\}^n\to[-1,1]$ depend on disjoint variable sets, and define $f(x)=\prod_{i=1}^kf_i(x)$. If every $\delta$-biased distribution $\epsilon$-fools every $f_i$, then every $\delta^k$-biased distribution fools $f$ with error $16^k\cdot2\epsilon$.
\end{authorclaim}
\subsection*{3.3 Limited Independence}
For $p\in[0,1]$, let $\operatorname{Bernoulli}(p)^{\otimes n}$ denote the distribution over $\{0,1\}^n$ where the bits are i.i.d. and each bit has expectation $p$. For example, $U_n=\operatorname{Bernoulli}(1/2)^{\otimes n}$. For a set $I=\{i_1<i_2<\cdots<i_\ell\}\subseteq[n]$ and a string $z\in\{0,1\}^n$, we let $z|_I=z_{i_1}z_{i_2}\cdots z_{i_\ell}\in\{0,1\}^\ell$. We say that $Z\in\{0,1\}^n$ is $\gamma$-almost $k$-wise independent with marginals $p$ if for every set $I\subseteq[n]$ with $|I|\le k$, the total variation distance between $Z|_I$ and $\operatorname{Bernoulli}(p)^{\otimes|I|}$ is at most $\gamma$.
\begin{authorclaim}{Claim7}
For every $n,k,C\in\mathbb N$ and $\gamma>0$, there is an explicit $\gamma$-almost $k$-wise independent distribution with marginals $p=1-2^{-C}$ that can be sampled with $O(Ck+\log(1/\gamma)+\log\log n)$ truly random bits.
\end{authorclaim}
\begin{proof}
Sample $Y\in\{0,1\}^{Cn}$ from a $(Ck)$-wise $(2^{-Ck/2-1}\gamma)$-biased distribution. Note that as discussed above $Y$ can be sampled using $O(\log(2^{Ck}/\gamma)+\log\log n)=O(Ck+\log(1/\gamma)+\log\log n)$ truly random bits. Divide $Y$ into $n$ blocks $Y^{(1)},\ldots,Y^{(n)}\in\{0,1\}^C$, and set $Z_i=0\Longleftrightarrow Y^{(i)}=1^C$. The desired distribution is $Z\in\{0,1\}^n$.
To prove correctness, let $f:\{0,1\}^n\to\{0,1\}$ be any test function depending on only $k$ variables. There is a function $g:\{0,1\}^{Cn}\to\{0,1\}$ depending on only $Ck$ variables such that $f(Z)=g(Y)$. By Claim5,
\[
|\mathbb E[f(Z)]-\mathbb E[f(\operatorname{Bernoulli}(p)^{\otimes n})]|
=|\mathbb E[g(Y)]-\mathbb E[g]|
\le2^{-Ck/2-1}\cdot2\gamma\cdot L_1(g)\le\gamma.
\]
\end{proof}
The expectation parameter $p$ can be "amplified" by drawing independent samples and combining with a coordinate-wise conjunction:
\begin{authorclaim}{Claim8}
Let $Z$ be $\gamma$-almost $k$-wise independent with marginals $p$. Draw $t$ independent samples $z^{(1)},\ldots,z^{(t)}\sim Z$, and let $Z'=z^{(1)}\wedge\cdots\wedge z^{(t)}$. Then $Z'$ is $(t\gamma)$-almost $k$-wise independent with marginals $p^t$.
\end{authorclaim}
\begin{proof}[Proof sketch]
The proof is a simple hybrid argument. Draw $t$ independent samples $r^{(1)},\ldots,r^{(t)}\sim\operatorname{Bernoulli}(p)^{\otimes n}$, and let $Z^{(i)}=z^{(1)}\wedge\cdots\wedge z^{(i)}\wedge r^{(i+1)}\wedge\cdots\wedge r^{(t)}$. One can show by induction on $i$ that $Z^{(i)}$ is $(i\gamma)$-almost $k$-wise independent with marginals $p^t$.
\end{proof}
\subsection*{3.4 PARITY $\circ$ AND Formulas}
Recall that our main result (Theorem1) is a PRG for read-once depth-2 $\mathrm{AC}^0[\oplus]$. For most of the paper, we will focus on the special case that the root gate is $\oplus$ and its immediate children are $\wedge$ gates. That is, define a PARITY $\circ$ AND formula to be a function of the form $f(x)=\bigoplus_{i=1}^mf_i(x)$, where each $f_i$ is a conjunction of literals, i.e., variables or their negations. We refer to $f_1,\ldots,f_m$ as the terms of $f$. We say that the formula is read-once if each variable $x_i$ appears in at most one term. Most of our effort will be spent fooling read-once PARITY $\circ$ AND formulas. Note that this is a slight generalization of read-once $\mathbb F_2$-polynomials due to the availability of $\neg$ gates. We will explain in Section5.6 why it is sufficient to focus on this special case.
The width of a term is the number of variables in the term; the width of $f$ is the maximum width of its terms. The length of $f$ is $m$, the number of its terms.
For convenience, if $f$ is a function taking values in $\{0,1\}$, we let $\mathfrak f=(-1)^f$. That way, if $f$ is a PARITY $\circ$ AND formula, $\mathfrak f=\prod_{i=1}^m\mathfrak f_i$.
\subsection*{3.5 Restrictions}
A restriction is a string $x\in\{0,1,?\}^n$; intuitively, $x_i=?$ means that $x_i$ has still not been assigned a value. We define an associative composition operation on restrictions by the formula $(x\circ x')_i=x_i$ if $x_i\ne?$, and $x'_i$ otherwise.
For a function $f$ on $\{0,1\}^n$, the restricted function $f|_x$ on $\{0,1\}^n$ is defined by $f|_x(x')=f(x\circ x')$.
A restriction $x$ can be specified by two strings $y,z\in\{0,1\}^n$ using the following notation.\footnote{With apologies, we here flip the order of the arguments to Res compared to the notation used in the authors' prior work [13].} Define $\operatorname{Res}:\{0,1\}^n\times\{0,1\}^n\to\{0,1,?\}^n$ by $(\operatorname{Res}(y,z))_i=?$ if $z_i=1$, and $y_i$ if $z_i=0$. In words, $z$ indicates the ? positions, and $y$ provides the bits in the non-? positions.
\subsection*{3.6 Pseudorandom Restrictions}
Let $Y,Z$ be distributions over $\{0,1\}^n$, and let $X=\operatorname{Res}(Y,Z)$. For a function $f:\{0,1\}^n\to\mathbb R$, we say that the distribution $X$ preserves the expectation of $f$ with error $\epsilon$ if $|\mathbb E[f|_X(U)]-\mathbb E[f]|\le\epsilon$.
An equivalent condition is that $|\mathbb E[f(Y+Z\wedge U)]-\mathbb E[f]|\le\epsilon$, where $+$ denotes addition over $\mathbb F_2^n$ and $\wedge$ denotes coordinate-wise conjunction. This second condition is the "pseudorandomness plus noise" perspective [18] (the string $Z\wedge U$ can be thought of as a noise vector.)
If $f$ takes on values in $\{0,1\}$, for each particular value $z$ that $Z$ might take on, we define the bias function [16] $\widetilde f_z:\{0,1\}^n\to[-1,1]$ by $\widetilde f_z(x)=\mathbb E[\mathfrak f(x+z\wedge U)]$. (We use $\mathfrak f$ rather than $f$ simply for convenience.) The statement that $X$ preserves the expectation of $f$ with error $\epsilon$ is also equivalent to the condition
\[
\left|\mathbb E_Z[\mathbb E_Y[\widetilde f_Z(Y)]-\mathbb E[\mathfrak f]]\right|\le2\epsilon.
\]
When $z$ is clear from context, we will just write $\widetilde f$ instead of $\widetilde f_z$.
If $X$ is a distribution over $\{0,1,?\}^n$ and $t\in\mathbb N$, let $X^{\circ t}$ denote the distribution over $x\in\{0,1,?\}^n$ obtained by drawing independent samples $x^{(1)},\ldots,x^{(t)}\sim X$ and composing them, $x=x^{(1)}\circ\cdots\circ x^{(t)}$.
Suppose $\mathcal F$ is a class of Boolean functions that is closed under restriction. If $X$ preserves the expectation of every $f\in\mathcal F$ with error $\epsilon$, then $X^{\circ t}$ preserves the expectation of every $f\in\mathcal F$ with error $t\epsilon$. Furthermore, informally, if $X$ "has ?-probability $p$", then $X^{\circ t}$ "has ?-probability $p^t$". To be precise, we can consider the case $X=\operatorname{Res}(Y,Z)$ where $Z$ is $\gamma$-almost $k$-wise independent with marginals $p$. Then the distribution of ? positions in $X^{\circ t}$ is described by Claim8.

\section*{4 Applying a Single Restriction}
In this section, we prove that the expectation of a PARITY $\circ$ AND formula is preserved under a suitable pseudorandom restriction. The cost of the restriction is only $O(\log n)$ truly random bits, the error is $\exp(-\widetilde\Omega(\log n))$ (near-optimal), and the restriction assigns values to a constant fraction of the inputs.
\subsection*{4.1 Restriction Construction}
Set $C=500$, $\bar C=2000C$, $c=1.1$, and $\beta=0.95$, and consider the following two distributions.
\begin{itemize}
\item Let $Y$ be a $\delta^3$-biased distribution over $\{0,1\}^n$ for $\delta=\min\{n^{-12\bar C},\frac12n^{-5c/(c-1)-1}\}=n^{-12,000,000}$.\footnote{No attempt was made to optimize the constants.}
\item Let $Z$ be a $\gamma$-almost $k$-wise independent distribution over $\{0,1\}^n$ with marginals $p=1-2^{-C}$, for $k=6\log n$ and $\gamma=n^{-9}$.
\end{itemize}
Our restriction is $\operatorname{Res}(Y,Z)$, i.e., $Z$ indicates where to put ? and $Y$ fills in the non-? bits.
\begin{authorclaim}{Lemma9}
Let $f$ be a read-once PARITY $\circ$ AND formula over $n$ variables of width at most $C\log n$. Then, $\operatorname{Res}(Y,Z)$ preserves the expectation of $f$ to within error $2^{-C\log n/\log\log n}$, i.e.,
\[
|\mathbb E[f(Y+Z\wedge U)]-\mathbb E[f]|\le2^{-C\log n/\log\log n}.
\]
\end{authorclaim}
\subsection*{4.2 Buckets}
Toward proving Lemma9, we first set some preliminary notations. Recall that $f$ is of the form $f=\bigoplus_{i=1}^mf_i=\bigoplus_{i=1}^m\bigwedge_{j=1}^{w_i}\ell_{ij}$, where every literal $\ell_{ij}$ is either some variable in $\{x_1,\ldots,x_n\}$ or its negation.
Set $Q=\lceil\log_c(C\log n)\rceil=O(\log\log n)$. We partition the terms of $f$ into $Q$ buckets according to their width. Namely, for each $q\in[Q]$ we define the interval $I_q=[c^{q-1},c^q)$ and define $B_q\subseteq[m]$ to be the set of indices $i$ such that $w_i\in I_q$. Also, for $q\in[Q]$ we define $F_q=\bigoplus_{i\in B_q}f_i$, so $f=\bigoplus_{q=1}^QF_q$. For every $q\in[Q]$ we further denote $m_q=|B_q|$.
We divide into two cases (Section4.3 and Section4.4) depending on whether there exists a bucket with substantially many terms. Lemma9 will follow immediately from Lemma10 and Lemma15, which cover these two cases respectively.
\subsection*{4.3 Case I -- There Exists a Heavy Bucket}
Say that bucket $q\in[Q]$ is heavy if both $m_q>3^{c^q}$ and $m_q>\log^C n$. The first case is that there exists a heavy bucket (i.e., there are many terms of roughly the same width, even relative to $q$). In this case, we will argue that $f$ itself is balanced and also that it stays balanced, w.h.p., after a pseudorandom restriction.
\begin{authorclaim}{Lemma10}
Let $f$ be a read-once PARITY $\circ$ AND formula over $n$ variables of width at most $C\log n$. Suppose there exists a heavy bucket as defined above. Then, with probability at least $1-1/n$ over $(y,z)\sim Y\times Z$,
\[
|\mathbb E[f|_{\operatorname{Res}(y,z)}]-\mathbb E[f]|\le1/n.
\]
\end{authorclaim}
Toward proving Lemma10, let us define a few more auxiliary notations. Write $f=f_{\mathrm{rest}}\oplus F_q$, where $q$ is a heavy bucket.
\begin{authorclaim}{Claim11}
It holds that $|\mathbb E[\mathfrak f]|\le1/(4n)$.
\end{authorclaim}
\begin{proof}
By the read-once property and the fact that $f_{\mathrm{rest}}$ is bounded,
\[
|\mathbb E[\mathfrak f]|=|\mathbb E[\mathfrak f_{\mathrm{rest}}]\mathbb E[\mathfrak F_q]|
\le|\mathbb E[\mathfrak F_q]|=\prod_{i\in B_q}|\mathbb E[\mathfrak f_i]|.
\]
Each term in $F_q$ has width at least $c^{q-1}$, so
\[
|\mathbb E[\mathfrak f]|\le(1-2\cdot2^{-c^{q-1}})^{m_q}\le e^{-2\cdot2^{-c^{q-1}}\cdot m_q}.
\]
Recalling that $m_q\ge3^{c^q}$, we have $2^{-c^{q-1}}\ge m_q^{-\gamma}$ for $\gamma=\log_32^{c^{-1}}<3/4$. Thus, using that fact that $m_q\ge\log^C n$,
\[
|\mathbb E[\mathfrak f]|\le e^{-2m_q^{1-\gamma}}\le e^{-2\log^{(1-\gamma)C}n}\le2^{-\log^{100}n}.
\]
\end{proof}
Next, we must analyze the bias of $f$ after the pseudorandom restriction. Let $n_q$ be the number of variables read by $F_q$. Let $b=\lceil\log_3n_q\rceil$. We will group the terms of $F_q$ into blocks, each of which reads roughly $b$ variables. To define this grouping, first observe that $b\ge\log_3m_q$, as each term reads at least one variable. Recalling that $c^q<\log_3m_q$, we know that $b>c^q$. Therefore, since each term in $F_q$ has width at most $c^q$, we can write $F_q=\bigoplus_{i=1}^Bg_i$, where each block $g_i$ reads $b_i$ variables for $b_i\in[b-\frac12c^q,b+\frac12c^q]$.
Let us now estimate $B$, the number of blocks. Since $b>c^q$, $b_i\in[b/2,3b/2]$. Also, $m_q>\log^C n$ so $b>(C/2)\log\log n$. Thus, on the one hand, $B\ge2n_q/(3b)\ge2\cdot3^b/(9b)$, and on the other hand, $B\le n_q\le3^b$.
Toward arguing that $f$ is balanced after pseudorandom restrictions, we wish to show that with high probability, $z\sim Z$ keeps many variables in many terms alive.
\begin{authorclaim}{Definition12}
For $z\in\{0,1\}^n$ and a formula $f$, we say $f$ is good under $z$ if $z$ assigns1 to at least a $(1-\beta)$-fraction of the variables $f$ reads.
\end{authorclaim}
\begin{authorclaim}{Claim13}
For a fixed $z\in\{0,1\}^n$, let $X_z\subseteq[B]$ be the set of blocks $g_i$ that are not good under $z$. Then, with probability at least $1-1/(2n)$ over $z\sim Z$, $|X_z|\le\lceil4\log n/b\rceil$.
\end{authorclaim}
\begin{proof}
Set $k_0=\lceil4\log n/b\rceil$. Let $S\subseteq[B]$ be some subset of cardinality $k_0$. We first bound the probability $p$ that every block $g_i$ for $i\in S$ is bad under $z\sim Z$. For a truly random $z\sim\operatorname{Bernoulli}(1-2^{-C})^{\otimes n}$, the above probability is bounded by
\[
\prod_{i\in S}\binom{b_i}{\beta b_i}2^{-C\beta b_i}
\le\prod_{i\in S}2^{b_i}2^{-5b_i}
\le(2^{-4\cdot b/2})^{|S|}\le n^{-8}.
\]
Now, for every $i\in[B]$, $k\ge k_0b_i$ so for $z\sim Z$, we get that $p\le n^{-8}+\gamma\le2n^{-8}$. Thus, by the union bound, with probability at most
\[
\binom B{k_0}p\le2B^{k_0}n^{-8}
\le2(3^b)^{4\log n/b}n^{-8}
\le2n^{-(2-\log3)4}\le n^{-4/3}<\frac1{2n}
\]
there will be some $S$ whose all blocks are bad. Taking the contrapositive, we infer that with probability at least $1-1/(2n)$ over $z\sim Z$, at most $k_0$ of the $g_i$-s are bad under $z$.
\end{proof}
\begin{authorclaim}{Lemma14}
With probability at least $1-1/n$ over $(y,z)\sim Y\times Z$, it holds that $|\mathbb E[\mathfrak f|_{\operatorname{Res}(y,z)}]|\le1/(2n)$.
\end{authorclaim}
\begin{proof}
Fix a good $z$, for which at most $4\log n/b$ of the $g_i$-s are not good under it. By Claim13, $z$ is good with probability at least $1-1/(2n)$. Let $B^{\mathrm{alive}}=[B]\setminus X_z$, so
\[
f=f_{\mathrm{rest}}\oplus\left(\bigoplus_{i\in B^{\mathrm{alive}}}g_i\right)\oplus\left(\bigoplus_{i\in[B]\setminus B^{\mathrm{alive}}}g_i\right).
\]
For every $i\in B^{\mathrm{alive}}$, set the following notations.
\begin{itemize}
\item For every $y\in\{0,1\}^n$, let $g_i^y$ denote the function $g_i|_{\operatorname{Res}(y,z)}$.
\item Let $I_i^{\mathrm{dead}}\subseteq[n]$ be the literals read by $g_i$ for which $z=0$. As $i\in B^{\mathrm{alive}}$, $|I_i^{\mathrm{dead}}|\le\beta b_i\le(3\beta/2)b$. Note that each literal $j\in I_i^{\mathrm{dead}}$ is set by $y\sim Y$.
\item Let $I_i^{\mathrm{alive}}\subseteq[n]$ be the literals read by $g_i$ for which $z=1$. As $i\in B^{\mathrm{alive}}$, $|I_i^{\mathrm{alive}}|\ge(1-\beta)b_i\ge((1-\beta)/2)b$.
\end{itemize}
Define the function $h_i$ so that $h_i(y)=1$ if $g_i^y$ is a nonconstant function, and0 otherwise. Namely, $h_i(y)=\bigwedge_{j\in I_i^{\mathrm{dead}}}y'_j$, where $y'_j$ is either $y_j$ or $\neg y_j$ depending on whether $y_j$ appears positively or negatively in $g_i$. Also, define $S(y)=\sum_{i\in B^{\mathrm{alive}}}h_i(y)$, where the sum is over the reals. Denote $\mu=\mathbb E[S(U)]=\sum_{i\in B^{\mathrm{alive}}}2^{-|I_i^{\mathrm{dead}}|}$, and note that $\mu\ge|B^{\mathrm{alive}}|\cdot2^{-3\beta b/2}$.
Set $\Delta S=S-\mu$. The spectral norm of the AND function is1, and so by the sub-additivity we get that $L_1(\Delta S)\le2|B^{\mathrm{alive}}|$. Set $\ell=2\lceil C\log n/(2\log(2|B^{\mathrm{alive}}|))\rceil$. By the sub-multiplicativity of the spectral norm we have that $L_1(\Delta S^\ell)\le(2|B^{\mathrm{alive}}|)^\ell\le n^C$.
For $\epsilon=1/2$, note that $\delta\le(\epsilon/2)\cdot L_1(\Delta S^\ell)^{-1}$. By Claim5, $Y$ $\epsilon$-fools the function $\Delta S^\ell$, so
\begin{equation}
|\mathbb E[(S(Y)-\mu)^\ell]-\mathbb E[(S(U)-\mu)^\ell]|\le\epsilon.\tag{5}
\end{equation}
Next, observe that $\Delta S(U)$ is the sum of zero-mean independent random variables, as the $h_i$-s are supported over disjoint set of variables. Set $A=|B^{\mathrm{alive}}|\cdot2^{-4\beta b}$. By the Chernoff bound,
\begin{align*}
\mathbb E[\Delta S(U)^\ell]
&\le\mathbb E[\Delta S(U)^\ell\mid\Delta S(U)^\ell\le A^\ell]\\
&\quad+\mathbb E[\Delta S(U)^\ell\mid\Delta S(U)^\ell\ge A^\ell]\cdot\Pr[\Delta S(U)^\ell\ge A^\ell]\\
&\le A^\ell+|B^{\mathrm{alive}}|^\ell\cdot\Pr[\Delta S(U)\ge A]\\
&\le|B^{\mathrm{alive}}|^\ell\cdot(2^{-4\beta b\ell}+e^{-2A^2/|B^{\mathrm{alive}}|}).
\end{align*}
Recall that $b>(C/2)\log\log n$, so $3^b\ge36\log n$ for a large enough $n$, and since $B\ge(2/(9b))3^b$ we get that $B\ge8\log n/b$ and $|B^{\mathrm{alive}}|\ge B-4\log n/b\ge B/2$. Next, we observe that
\[
\frac{2A^2}{|B^{\mathrm{alive}}|}=2|B^{\mathrm{alive}}|2^{-8\beta b}
\ge B\cdot2^{-8\beta b}
\ge\frac2{9b}2^{(\log3-8\beta)b}\ge2^b.
\]
As $b\ell\le Cb\log n/\log B\le C\log n$, we can conclude that $2^b\ge4\beta b\ell$ and so $e^{-2A^2/|B^{\mathrm{alive}}|}\le2^{-4\beta b\ell}$, which implies that $\mathbb E[\Delta S(U)^\ell]\le2|B^{\mathrm{alive}}|^\ell\cdot2^{-4\beta b\ell}$.
Using Equation(5) and the above bound yields a bound on $\mathbb E[\Delta S(Y)^\ell]$. By Markov's inequality,
\begin{equation}
\Pr[S(Y)<\mu/2]\le\frac{\mathbb E[(S(Y)-\mu)^\ell]}{(\mu/2)^\ell}
\le\frac{\epsilon+2|B^{\mathrm{alive}}|^\ell\cdot2^{-4\beta b\ell}}{(\mu/2)^\ell}
\le\left(\frac{8|B^{\mathrm{alive}}|2^{-4\beta b}}{\mu}\right)^\ell.\tag{6}
\end{equation}
Recalling that $\mu\ge|B^{\mathrm{alive}}|\cdot2^{-3\beta b/2}$, Equation(6) becomes
\[
\Pr[S(Y)<\mu/2]\le(8\cdot2^{(-4\beta+3\beta/2)b})^\ell
<2^{-2\beta b\ell}\le2^{-\beta C\log n/2}\le\frac1{2n},
\]
where we have used the fact that $b\ell\ge C\log n/4$.
Overall, with probability at least $1-1/(2n)$ over $y\sim Y$, $g_i^y$ is nonconstant for at least $\mu/2$ of the $i$-s, and recall that each such $g_i^y$ is over at least $(1-\beta)b_i$ variables. Fix such a good $y$, and let $\mathcal G\subseteq[B^{\mathrm{alive}}]$ be the set of nonconstant $g_i^y$-s. Again, we can write
\[
\bigoplus_{i\in B^{\mathrm{alive}}}g_i^y
=\left(\bigoplus_{i\in\mathcal G}g_i^y\right)\oplus\left(\bigoplus_{i\in B^{\mathrm{alive}}\setminus\mathcal G}g_i^y\right)
\mathrel{\mathop:}=t_1\oplus t_2.
\]
Similarly to Claim11, in order to bound the bias of $f|_{\operatorname{Res}(Y,Z)}$ it is sufficient to bound the bias of $t_1$, and so
\[
|\mathbb E[\mathfrak t_1]|\le(1-2^{-3b/2})^{\mu/2}.
\]
Using the fact that $\mu\ge(B/2)\cdot2^{-3\beta b/2}\ge(1/(9b))2^{(\log3-3\beta/2)b}>2^{301b/200}$, we get
\[
|\mathbb E[\mathfrak t_1]|\le e^{-2^{-3b/2}2^{301b/200}}\le e^{-\log^{C/400}n}\le\frac1{2n}.
\]
\end{proof}
\begin{proof}[Proof of Lemma10]
Finally, the fact that with probability at least $1-1/n$ over $(y,z)\sim Y\times Z$, $|\widetilde f_z(y)-\mathbb E[\mathfrak f]|\le1/n$, follows immediately from Claim11 and Lemma14.
\end{proof}

\subsection*{4.4 Case II -- There Are No Heavy Buckets}
In this subsection, we prove that a single pseudorandom restriction preserves the expectation in the case where there is no such a heavy $B_q$. Namely, for every $q\in[Q]$, either $m_q\le3^{c^q}$ or $m_q\le\log^C n$ (or both).
\begin{authorclaim}{Lemma15}
Let $f$ be a read-once PARITY $\circ$ AND formula over $n$ variables in which the width of every term is at most $C\log n$, and in which there are no heavy buckets as described above. Then, with probability at least $1-\frac12\cdot2^{-C\log n/\log\log n}$ over $z\sim Z$ it holds that
\[
|\mathbb E[\widetilde f_z(Y)]-\mathbb E[\mathfrak f]|
\le\frac12\cdot2^{-C\log n/\log\log n}.
\]
\end{authorclaim}
Toward proving Lemma15, we partition the $Q$ buckets into two sets and treat terms that fall into each set of buckets separately. Namely, define the two sets as follows. $\mathcal A=\{q\in[Q]:m_q\le\log^{2C}n\}$. We refer to these buckets as the sparse buckets. $\mathcal B=[Q]\setminus\mathcal A$. We refer to these buckets as the well-behaved buckets.
For each set $T\in\{\mathcal A,\mathcal B\}$ we denote $f_T=\bigoplus_{i\in T}F_i$, and so $f=f_{\mathcal A}\oplus f_{\mathcal B}$. The next two subsections will be devoted to proving that the expectation of each $f_T$ is preserved after a single pseudorandom restriction. In Section4.4.3 we will combine the two results using the XOR lemma for small-bias distributions (Lemma6) to prove Lemma15.
\subsection*{4.4.1 Handling Sparse Buckets}
For the sparse buckets, we will follow the Forbes-Kelley approach [14] to prove the following.
\begin{authorclaim}{Lemma16}
With probability at least $1-\frac14\cdot2^{-C\log n/\log\log n}$ over $z\sim Z$, it holds that
\[
|\mathbb E[(\widetilde{f_{\mathcal A}})_z(Y)]-\mathbb E[\mathfrak f_{\mathcal A}]|
\le\frac14\cdot2^{-C\log n/\log\log n}.
\]
\end{authorclaim}
As outlined in Section1.5.1, Lemma16 follows readily from the work by Forbes and Kelley [14]. We require our restriction to work with high probability over $z\sim Z$, not merely in expectation, so we must redo some of Forbes and Kelley's analysis. (No substantial modification is needed.) The details follow.
\begin{proof}[Proof of Lemma16]
First, recall that each term in $f_{\mathcal A}$ is of width at most $C\log n$. There are at most $\log^{2C}n$ terms in each bucket, and at most $Q=O(\log\log n)$ such buckets, so overall $f_{\mathcal A}$ reads at most $n'=\log^{2C+2}n$ variables.
Note that $f_{\mathcal A}$ can be computed by a width-4 ROBP of length $n'$. We follow [14] and let $G:\{0,1\}^{n'}\to\mathbb R^{4\times4}$ encode the transition of the branching program. Namely, perhaps after renumbering the variables, we have $G(x)=G_1(x_1)\cdots G_{n'}(x_{n'})$ where $G_i(x_i)=A_{i,x_i}$ for $A_{i,b}$ being the transition matrix that corresponds to taking the bit $b$ while at layer $i$. Set $k_0=8\log n/\log\log n$, and note that $k_0\le k$. By [14, Lemma4.1], $G$ can be written as
\[
G=\mathbb E[G]+L+\sum_{i=1}^{n'}H_i\cdot G^{>i},
\]
where $L$ has degree\footnote{We say a function $H:\{0,1\}^n\to\mathbb R^{w\times w}$ having Fourier expansion $\sum_{\alpha\in\mathbb F_2^n}\widehat H_\alpha\chi_\alpha$ has degree $d$ if $\widehat H_\alpha$ is the zero matrix for every $\alpha$ with Hamming weight larger than $d$.} less than $k_0$, $H_i$ is of degree exactly $k_0$, $G^{>i}$ is a width-4 ROBP, and $H_i$ and $G^{>i}$ are on disjoint set of variables. More specifically,
\[
L=\sum_{\substack{\alpha\in\mathbb F_2^{n'}\\0<|\alpha|<k_0}}\widehat G_\alpha\chi_\alpha
\]
is the truncated Fourier expansion of $G$, $G^{>i}(x_{i+1},\ldots,x_n)=G_{i+1}(x_{i+1})\cdots G_{n'}(x_{n'})$, and
\[
H_i=\sum_{\substack{\alpha\in\mathbb F_2^{n'}\\|\alpha|=k_0,\alpha_i=1}}\widehat{G^{\le i}}_\alpha\chi_\alpha,
\]
where $G^{\le i}(x_1,\ldots,x_i)=G_1(x_1)\cdots G_i(x_i)$. Let $\|\cdot\|$ be the Frobenius norm. By sub-additivity, we have
\begin{align}
\mathbb E_Z[\|\mathbb E_{Y,U}[G(Y+Z\wedge U)]-\mathbb E[G]\|]
&\le\mathbb E_Z[\|\mathbb E_{Y,U}[L(Y+Z\wedge U)]\|]\notag\\
&\quad+\sum_{i=1}^{n'}\mathbb E_Z[\|\mathbb E_{Y,U}[(H_i\cdot G^{>i})(Y+Z\wedge U)]\|].\tag{7}
\end{align}
Just as in [14], the low-degree term $L$ is dealt with a $\delta$-biased distribution. From the work of Chattopadhyay, Hatami, Reingold, and Tal [9] we know that
\[
L_1(L)=\sum_{k'_0=1}^{k_0}(c_{\mathrm{CHRT}}\log n')^{4k'_0}\le2(c_{\mathrm{CHRT}}\log n')^{4k_0}
\]
for some universal constant $c_{\mathrm{CHRT}}\ge1$. Thus, by Claim5, we get that the first term of Equation(7) is bounded by
\[
2\delta\cdot2(c_{\mathrm{CHRT}}\log n')^{4k_0}
\le2^{-C\log n}\cdot2^{8k_0\log\log\log n}\le n^{-C/2},
\]
taking into account the fact that $\mathbb E[L(U)]=0$.
For each $i$ of the second term of Equation(7), we use sub-multiplicativity and the fact that $H_i$ and $G^{>i}$ are on disjoint set of variables to get
\[
\mathbb E_Z[\|\mathbb E_{Y,U}[(H_i\cdot G^{>i})(Y+Z\wedge U)]\|]
\le\mathbb E_{Y,Z}[\|\mathbb E_U[H_i(Y+Z\wedge U)]\|\cdot\|\mathbb E_U[G^{>i}(Y+Z\wedge U)]\|].
\]
As $G^{>i}$ is a width-4 ROBP, $\|\mathbb E_U[G^{>i}(y+z\wedge U)]\|\le2$ for all $y\sim Y$ and $z\sim Z$. Continuing the above bound, by Cauchy-Schwarz we get
\[
\mathbb E_Z[\|\mathbb E_{Y,U}[(H_i\cdot G^{>i})[Y+Z\wedge U]]\|]
\le2\sqrt{\mathbb E_{Y,Z}[\|\mathbb E_U[H_i(Y+Z\wedge U)]\|^2]}.
\]
Following [14, Lemma7.1]\footnote{Forbes and Kelley [14] take the bits of $Z$ to have marginals $p=1/2$, but one can extend the lemma easily for the case of a general $p$.}, using the bound by Chattopadhyay et al. [9] and Parseval's identity [14, Proposition3.1], we get
\begin{align*}
\mathbb E_{Y,Z}[\|\mathbb E_U[H_i(Y+Z\wedge U)]\|^2]
&\le(2^{-Ck_0}+\gamma)\cdot\left(\delta\left(\sum_{\alpha\in\mathbb F_2^{n'}}\|(\widehat H_i)_\alpha\|\right)^2+\sum_{\alpha\in\mathbb F_2^{n'}}\|(\widehat H_i)_\alpha\|^2\right)\\
&\le(2^{-Ck_0}+\gamma)\cdot\left(\delta\cdot L_1^2(G^{\le i})+\mathbb E[\|G^{\le i}(U)\|^2]\right)
\le8\cdot2^{-Ck_0}.
\end{align*}
Overall, we get that
\begin{align*}
\mathbb E_Z[\|\mathbb E_{Y,U}[G(Y+Z\wedge U)]-\mathbb E[G]\|]
&\le n^{-C/2}+2n'\sqrt{8\cdot2^{-Ck_0}}\\
&\le\frac1{16}\cdot2^{-Ck_0/4}
=\frac1{16}\cdot2^{-2C\log n/\log\log n},
\end{align*}
and we can choose the encoding $G$ so that $f_{\mathcal A}(x)=G(x)_{1,1}$. Markov's inequality completes the proof.
\end{proof}

\subsection*{4.4.2 Handling Well-Behaved Buckets}
We will use our tail bounds for subset-wise symmetric polynomials to prove the following lemma.
\begin{authorclaim}{Lemma17}
With probability at least $1-1/(2n)$ over $z\sim Z$, $f_{\mathcal B}$ can be written as $f_{\mathcal B}=f'_{\mathcal B}\oplus f''_{\mathcal B}$, where $f'_{\mathcal B}$ and $f''_{\mathcal B}$ are over disjoint set of variables, and for every $g\in\{f'_{\mathcal B},f''_{\mathcal B}\}$ it holds that $|\mathbb E[\widetilde g_z(Y)]-\mathbb E[\mathfrak g]|\le1/n$.
\end{authorclaim}
The proof of Lemma17 will follow immediately from Claim20 and Lemma21. Toward proving the above lemma, let us set some preliminaries.
\begin{authorclaim}{Claim18}
If $q\in\mathcal B$ then $c^q\in[C\log\log n,C\log n]$ and $m_q\le3^{c^q}$.
\end{authorclaim}
\begin{proof}
The upper bound on $c^q$ follows immediately from the assumption in Lemma15 that every term has width at most $C\log n$. Also, $m_q>\log^{2C}n$ since $q\notin\mathcal A$. Since we are at CaseII, $m_q>\log^{2C}n$ implies that $m_q\le3^{c^q}$. From the fact that $\log^{2C}n<3^{c^q}$ we get $c^q>\log_3(\log^{2C}n)>C\log\log n$.
\end{proof}
Recall that a term $f_i$ is good under $z$ if the variables read by $f_i$ intersects with $z$ in at least $1-\beta$ fraction.
\begin{authorclaim}{Claim19}
For a fixed $z\in\{0,1\}^n$, let $X_z\subseteq[m]$ be the set of terms in $f_{\mathcal B}$ that are not good under $z$. Then, with probability at least $1-1/(2n)$ over $z\sim Z$, $|X_z|\le(3c/(c-1))\log n$.
\end{authorclaim}
\begin{proof}
The proof is very similar to Claim13. Fix a bucket $q\in\mathcal B$, set $k_q=3\log n/c^q$ and observe that $k\ge k_q$. Let $S\subseteq B_q$ be some subset of cardinally $k_q$. We first bound the probability $p$ that every term $f_i$ for $i\in S$ is bad under $z\sim Z$.
For a truly random $z\sim\operatorname{Bernoulli}(1-2^{-C})^{\otimes n}$, the above probability is bounded by
\[
\prod_{i\in S}\binom{w_i}{\beta w_i}2^{-\beta Cw_i}
\le\prod_{i\in S}2^{w_i}2^{-5w_i}
\le(2^{-4c^{q-1}})^{k_q}
\le2^{-3k_qc^q}\le n^{-9}.
\]
For $z\sim Z$, we get that $p\le n^{-9}+\gamma\le2n^{-9}$. Thus, with probability at most $\binom{m_q}{k_q}p$ over $z\sim Z$ there exists a set of $k_q$ terms in $B_q$ whose all terms are bad under $z$. By using Claim18, we get
\[
\binom{m_q}{k_q}p\le m_q^{k_q}\cdot2^{-9\log n+1}
\le3^{k_qc^q+\log_32\cdot(-9\log n+1)}
\le3^{-9\log n/4}\le n^{-3}.
\]
Moreover, with probability at most $|\mathcal B|n^{-3}\le n^{-2}$ over $z\sim Z$ there exists a $q\in\mathcal B$ and a set of $k_q$ terms in $B_q$ whose all terms are bad under $z$. Taking the contrapositive, we infer that with probability at least $1-n^{-2}\ge1-1/(2n)$ over $z\sim Z$, we have at most
\[
\sum_{q\in\mathcal B}k_q\le\sum_{q=1}^Q\frac{3\log n}{c^q}\le\frac{3c}{c-1}\log n
\]
terms that are bad for $z$.
\end{proof}
From here onwards, we fix a $z$ satisfying $|X_z|\le(3c/(c-1))\log n$. Write
\[
f_{\mathcal B}=\bigoplus_{i\in\mathcal C\setminus X_z}f_i\oplus\bigoplus_{i\in X_z}f_i
\mathrel{\mathop:}=f'_{\mathcal B}\oplus f''_{\mathcal B},
\]
where $\mathcal C=\bigcup_{q\in\mathcal B}B_q\subseteq[m]$ is the set of all terms that belong to $\mathcal B$'s buckets. Simply put, we divide $f_{\mathcal B}$ to the parity of exceptional terms $f''_{\mathcal B}$ and non-exceptional terms $f'_{\mathcal B}$ for whom we will refer to as good terms. We stress that both $f'_{\mathcal B}$ and $f''_{\mathcal B}$ depend on $z$.
\begin{authorclaim}{Claim20 (Exceptional terms)}
$|\mathbb E[(\widetilde{f''_{\mathcal B}})_z(Y)]-\mathbb E[\mathfrak f''_{\mathcal B}]|\le1/n$.
\end{authorclaim}
\begin{proof}
For brevity, let $g=f''_{\mathcal B}$. For a fixed $w\in\{0,1\}^n$, let $g^w(x)=\mathfrak g(x+w)$. The proof will follow from bounding the spectral norm of $g^w$. Indeed, $g^w$ is a multiplication of at most $(3c/(c-1))\log n$ terms, each of which has spectral norm at most3. By sub-multiplicativity,
\[
L_1(g^w)\le3^{(3c/(c-1))\log n}\le n^{5c/(c-1)}.
\]
Now, $\delta\le\frac12n^{-5c/(c-1)-1}$, so by Claim5 we get that $|\mathbb E[g^w(Y)]-\mathbb E[g^w]|\le1/n$ for every $w\in\{0,1\}^n$.
Fooling $g^w$ is sufficient to fool $\widetilde g_z$. To see this, note that
\begin{align*}
|\mathbb E[\widetilde g_z(Y)]-\mathbb E[\mathfrak g]|
&=|\mathbb E[\mathfrak g(Y+z\wedge U)]-\mathbb E[\mathfrak g(U+z\wedge U')]|\\
&=\left|\mathbb E_{w\sim U}[\mathbb E[g^w(Y)]-\mathbb E[g^w]]\right|\le1/n,
\end{align*}
where $U'$ is an independent copy of $U$.
\end{proof}
Next, we prove:
\begin{authorclaim}{Lemma21 (Good terms)}
$|\mathbb E[(\widetilde{f'_{\mathcal B}})_z(Y)]-\mathbb E[\mathfrak f'_{\mathcal B}]|\le1/n$.
\end{authorclaim}
\begin{proof}
For brevity, let $g=f'_{\mathcal B}$ and recall that its set of terms is given by $\mathcal C\setminus X_z$. Shifting the bias function $\widetilde g=\widetilde g_z$ to mean zero, recall that we define $\check g(x)=\widetilde g(x)/\mathbb E[\widetilde g]-1$.
Thus, we can write
\begin{align*}
\widetilde g
&=\mathbb E[\mathfrak g]\prod_{i\in\mathcal C\setminus X_z}(1+\widehat g_i)
=\mathbb E[\mathfrak g]\sum_{I\subseteq\mathcal C\setminus X_z}\prod_{i\in I}\widehat g_i\\
&=\mathbb E[\mathfrak g]\sum_{\vec k\in\mathbb N^Q}
\sum_{\substack{I\subseteq\mathcal C\setminus X_z\\K(I)=\vec k}}\prod_{i\in I}\widehat g_i,
\end{align*}
where by $K(I)=\vec k$ we mean that for every $q\in[Q]$, there are $\vec k[q]$ terms in $I$ that belong to the $q$-th bucket, i.e., $|I\cap B_q|=\vec k[q]$. For simplicity, we reorder the terms of $g$ and write $g=\bigoplus_{i\in[m']}g_i$ for $m'=|\mathcal C\setminus X_z|$, and for $q\in[Q]$, $B_q\subseteq[m']$ is the set of terms in $g$ that belong to the $q$-th bucket. We abbreviate $\vec g=(\widehat g_1,\ldots,\widehat g_{m'})$, and write
\[
S_{\vec k}(\vec g)=\sum_{\substack{I\subseteq[m']\\K(I)=\vec k}}\prod_{i\in I}\widehat g_i.
\]
Under these notations, $\widetilde g=\mathbb E[\mathfrak g]\sum_{\vec k\in\mathbb N^Q}S_{\vec k}(\vec g)$.
Let $I^g(x)$ be the Boolean-valued function which is1 if and only if
\[
\sum_{\substack{\vec k\in\mathbb N^Q\\\|\vec k\|_{(c)}>A}}|S_{\vec k}(\vec g(x))|\le2^{-A/1024},
\]
where $A=\bar C\log n$ and $\|\vec k\|_{(c)}=\sum_{q=1}^Qc^q\vec k[q]$. Section4.4.4 will be devoted to showing that $\mathbb E[I^g(Y)]$ is very close to1. Namely,
\begin{authorclaim}{Lemma22}
The following two inequalities hold.
\begin{enumerate}
\item $\mathbb E[I^g(Y)]\ge1-e^{-c_IA}$ for $c_I=\ln2/2^{23}$.
\item $\mathbb E[S_{\vec k}^2(\vec g(Y))]\le2^{-\|\vec k\|_{(c)}/8}$.
\end{enumerate}
\end{authorclaim}
For now, let us take Lemma22 as given and continue with the proof of Lemma21. We proceed by writing
\begin{equation}
|\mathbb E[\widetilde g(Y)]-\mathbb E[\widetilde g]|
\le|\mathbb E[\widetilde g(Y)\mid I^g(Y)=1]-\mathbb E[\widetilde g]|
+2\Pr[I^g(Y)=0].\tag{8}
\end{equation}
By Lemma22, we have that $\Pr[I^g(Y)=0]\le e^{-c_IA}$. Next, observe that
\[
|\mathbb E[\widetilde g(Y)\mid I^g(Y)=1]-\mathbb E[\widetilde g]|
=\left|\mathbb E[\mathfrak g]\sum_{\substack{\vec k\in\mathbb N^Q\\\|\vec k\|_{(c)}>0}}
\mathbb E[S_{\vec k}(\vec g(Y))\mid I^g(Y)=1]\right|,
\]
and set
\[
\Delta=\left|\sum_{\substack{\vec k\in\mathbb N^Q\\\|\vec k\|_{(c)}>0}}
\mathbb E[S_{\vec k}(\vec g(Y))\mid I^g(Y)=1]\right|,
\]
so Equation(8) gives us
\begin{equation}
|\mathbb E[\widetilde g(Y)]-\mathbb E[\widetilde g]|\le\Delta+2e^{-c_IA}.\tag{9}
\end{equation}
We bound $\Delta$ as follows.
\begin{align*}
\Delta&\le\left|\sum_{\substack{\vec k\in\mathbb N^Q\\0<\|\vec k\|_{(c)}\le A}}
\mathbb E[S_{\vec k}(\vec g(Y))\mid I^g(Y)=1]\right|\\
&\quad+\max_{\substack{y\in\{0,1\}^n\\I^g(y)=1}}
\sum_{\substack{\vec k\in\mathbb N^Q\\\|\vec k\|_{(c)}>A}}|S_{\vec k}(\vec g(y))|.
\end{align*}
By definition, the second term is at most $2^{-A/1024}$. The first term, call it $\Delta_1$, can be split into two terms as follows.
\begin{align}
\Delta_1
&=\frac1{\Pr[I^g(Y)=1]}\left|\sum_{\substack{\vec k\in\mathbb N^Q\\0<\|\vec k\|_{(c)}\le A}}\mathbb E[S_{\vec k}(\vec g(Y))\cdot I^g(Y)]\right|\notag\\
&\le2\left|\sum_{\substack{\vec k\in\mathbb N^Q\\0<\|\vec k\|_{(c)}\le A}}\mathbb E[S_{\vec k}(\vec g(Y))\cdot I^g(Y)]\right|\tag{10}\\
&\le2\left|\sum_{\substack{\vec k\in\mathbb N^Q\\0<\|\vec k\|_{(c)}\le A}}\mathbb E[S_{\vec k}(\vec g(Y))]\right|
+2\left|\sum_{\substack{\vec k\in\mathbb N^Q\\0<\|\vec k\|_{(c)}\le A}}\mathbb E[S_{\vec k}(\vec g(Y))\cdot(1-I^g(Y))]\right|\notag\\
&\le2\left|\sum_{\substack{\vec k\in\mathbb N^Q\\0<\|\vec k\|_{(c)}\le A}}\mathbb E[S_{\vec k}(\vec g(Y))]\right|
+2\sqrt{\mathbb E[1-I^g(Y)]}\sum_{\substack{\vec k\in\mathbb N^Q\\0<\|\vec k\|_{(c)}\le A}}\sqrt{\mathbb E[S_{\vec k}^2(\vec g(Y))]}.\tag{11}
\end{align}
where the last inequality follows from the triangle inequality followed by Cauchy-Schwarz.
By Lemma22, the second term of Equation(11), $\Delta_{1,2}$, is at most
\begin{align*}
\Delta_{1,2}
&\le2e^{-c_IA}\sum_{\substack{\vec k\in\mathbb N^Q\\0<\|\vec k\|_{(c)}\le A}}\sqrt{2^{-\|\vec k\|_{(c)}/8}}\\
&\le2e^{-c_IA}\sum_{w=1}^{A-1}|\{\vec k\in\mathbb N^Q:w<\|\vec k\|_{(c)}\le w+1\}|2^{-w/\sqrt8}\\
&\le2e^{-c_IA}(A+1)^Q\sum_{w=1}^A2^{-w/\sqrt8}
\le8(A+1)^Qe^{-c_IA}\\
&\le2^{(2/\log c)(\log\log n)^2}e^{-c_IA}\le\frac1{8n}.
\end{align*}
To finish bounding $\Delta_1$, it is left to bound the first term of Equation(11), denoted by $\Delta_{1,1}$.
\begin{authorclaim}{Claim23}
$\Delta_{1,1}=2|\sum_{\vec k\in\mathbb N^Q,\ 0<\|\vec k\|_{(c)}\le A}\mathbb E[S_{\vec k}(\vec g(Y))]|\le1/(8n)$.
\end{authorclaim}
\begin{proof}
The proof goes by bounding the spectral norm of the function $S_{\vec k}(\vec g(x))$. As for every $\vec k\in\mathbb N^Q$ with $\|\vec k\|_{(c)}\ne0$, $\mathbb E[S_{\vec k}(\vec g(U))]=0$, the claim will follow by using Claim5, together with sub-additivity and sub-multiplicativity. First, note that:
\begin{authorclaim}{Claim24}
For every $i\in[m]$, $L_1(\widehat g_i)\le4$.
\end{authorclaim}
\begin{proof}
Consider the function $h_i=1-g_i$, so $L_1(\widetilde g_i)\le L_1(\widetilde h_i)+1$ and $\mathbb E[\widetilde h_i]=\mathbb E[h_i]=1-\mathbb E[g_i]\ge1/2$. Now,
$L_1(\widehat h_i)\le(1/\mathbb E[\widetilde h_i])L_1(\widetilde h_i)+1\le2L_1(\widetilde h_i)+1$. Recalling that $\widetilde h_i(x)=\mathbb E[h_i(x+z\wedge U)]$, we get $L_1(\widetilde h_i)\le1$ as every shift of $h_i$ is a negated conjunction of literals. Thus, $L_1(\widehat h_i)\le3$ and $L_1(\widehat g_i)\le4$.
\end{proof}
Then, for every such $\vec k\in\mathbb N^Q$,
\begin{align*}
L_1(S_{\vec k}(\vec g))
&\le\sum_{\substack{I\subseteq\mathcal C\setminus X_T\\K(I)=\vec k}}\prod_{i\in I}L_1(\widehat g_i)
\le\sum_{\substack{I\subseteq\mathcal C\setminus X_T\\K(I)=\vec k}}4^{|I|}\\
&=\sum_{\substack{I\subseteq\mathcal C\setminus X_T\\K(I)=\vec k}}\prod_{q\in[Q]}4^{\vec k[q]}
=\prod_{q\in[Q]}\binom{\vec k[q]}{m_q}4^{\vec k[q]}.
\end{align*}
Recall that Claim18 tells us that $m_q\le3^{c^q}$, so
\begin{equation}
L_1(S_{\vec k}(\vec g))\le\prod_{q\in[Q]}3^{c^q(1+\log_34)\vec k[q]}\le12^{\|\vec k\|_{(c)}}.\tag{12}
\end{equation}
Finally,
\[
L_1\left(\sum_{\substack{\vec k\in\mathbb N^Q\\0<\|\vec k\|_{(c)}\le A}}\mathbb E[S_{\vec k}(\vec g(Y))]\right)
\le(A+1)^Q\cdot12^A\le2^{6A}\le n^{6\bar C},
\]
and the claim follows by observing that $\delta\le(1/(32n))n^{-6\bar C}$.
\end{proof}
Incorporating the above claim, we get that $\Delta_1=\Delta_{1,1}+\Delta_{1,2}\le1/(8n)+1/(8n)\le1/(4n)$, which readily gives $\Delta\le1/(4n)+2^{-A/1024}\le1/(2n)$. Plugging-it in Equation(9), we finally get
\[
|\mathbb E[\widetilde g(Y)]-\mathbb E[\widetilde g]|
\le\frac1{2n}+2e^{-c_IA}\le\frac1n
\]
and the desired result.
\end{proof}

\subsection*{4.4.3 Putting It Together}
Here we finally incorporate Lemma16 and Lemma17.
\begin{proof}[Proof of Lemma15]
By Lemma16 and Lemma17, with probability at least $1-\frac14\cdot2^{-C\log n/\log\log n}-1/n\ge1-\frac12\cdot2^{-C\log n/\log\log n}$ over $z\sim Z$, we can write $f=f_{\mathcal A}\oplus f'_{\mathcal B}\oplus f''_{\mathcal B}$, where the three functions are over disjoint set of variables, and it holds that for each $T\in\{\mathcal A,\mathcal B,\mathcal B'\}$,
\[
|(\widetilde f_T)_z(Y')-\mathbb E[\mathfrak f_T]|
\le\frac14\cdot2^{-C\log n/\log\log n}
\]
for any $\delta$-biased distribution $Y'$. Using the XOR lemma for small-biased spaces (see Lemma6), taking into account that our distribution $Y$ is in fact $\delta^3$-biased, we conclude that
\[
|\mathbb E[\widetilde f_z(Y)]-\mathbb E[\mathfrak f]|
\le16^3\cdot2\cdot\frac14\cdot2^{-C\log n/\log\log n}
\le\frac12\cdot2^{-C\log n/\log\log n},
\]
and the lemma follows.
\end{proof}
\subsection*{4.4.4 $I^g$ Almost Always Happens}
We keep using the notations of Section4.4.2. Specifically, recall that $g=f'_{\mathcal B}=\bigoplus_{i\in[m']}g_i$ for $m'=|\mathcal C\setminus X_z|$, and for $q\in[Q]$, $B_q\subseteq[m']$ is the set of terms in $g$ that belong to the $q$-th bucket. Also, for $\vec g=(\widehat g_1,\ldots,\widehat g_{m'})$,
\[
S_{\vec k}(\vec g)=\sum_{\substack{I\subseteq[m']\\K(I)=\vec k}}\prod_{i\in I}\widehat g_i.
\]
Recall that $I^g(x)\in\{0,1\}$ is1 if and only if
\[
\sum_{\substack{\vec k\in\mathbb N^Q\\\|\vec k\|_{(c)}>A}}|S_{\vec k}(\vec g(x))|\le2^{-A/1024},
\]
where $A=\bar C\log n$ and $\|\vec k\|_{(c)}=\sum_{q\in[Q]}c^q\vec k[q]$.
\begin{proof}[Proof of Lemma22]
As in Section2, we define $R_{\vec k}(\vec g)=S_{\vec k}^2(\vec g)\cdot\prod_{q\in[Q]}\vec k[q]!$.
By Lemma3, to prove the bound on $\Pr[I^g(Y)=0]$ it is sufficient to prove that for every $\vec k\in\mathbb N^Q$ with $\|\vec k\|_{(c)}\le A$ we have that
\[
\mathbb E[R_{\vec k}(\vec g(Y))]\le2^{-\|\vec k\|_{(c)}/8}.
\]
By now a standard course of action, we aim at bounding the spectral norm of the function $R_{\vec k}(\vec g)$, together with its expectation under the uniform distribution. To this end, let us define, for $q\in[Q]$ and an integer $\ell$,
\[
\check S_{\ell,q}=\sum_{\substack{I\subseteq B_q\\|I|=\ell}}\prod_{i\in I}\widehat g_i,
\]
so $R_{\vec k}(\vec g)=\prod_{q\in[Q]}\check S_{\vec k[q],q}^2\vec k[q]!$. First, we record that:
\begin{authorclaim}{Claim25}
For every $i\in[m']$, $\mathbb E[\widehat g_i^2]\le2^{-(2-2\beta)w_i}$.
\end{authorclaim}
\begin{proof}
Let $V_i\subseteq[n]$ be the set of variables read by $g_i$, of cardinality $w_i$, and let $\ell_i=|V_i\cap\{j\in[n]:z_j=1\}|$ be the number of live variables read by $g_i$. Note that
\[
\widetilde g_i(x)=\mathbb E[g_i(x+z\wedge U)]
=\begin{cases}0&\text{if there exists }j\in V_i\text{ such that }x_j=z_j=0,\\2^{-\ell_i}&\text{otherwise}.\end{cases}
\]
Then,
\begin{align*}
\mathbb E[\widetilde g_i^2]
&=2^{-2\ell_i}\Pr_{x\sim U}[\text{for every }j\in V_i\text{ s.t. }z_j=0\text{ it holds that }x_j=1]\\
&=2^{-2\ell_i}2^{-(w_i-\ell_i)}=2^{-w_i-\ell_i}.
\end{align*}
Recalling that $\ell_i\ge(1-\beta)w_i$ ($g_i$ is good under $z$), we have $\mathbb E[\widetilde g_i^2]\le2^{-(2-\beta)w_i}$. Let $h_i=1-g_i$, and note that
\[
\mathbb E[\widehat h_i^2]=\operatorname{Var}[\widehat h_i]
=\frac{\operatorname{Var}[\widetilde g_i]}{\mathbb E[h_i]^2}
\le4\cdot\mathbb E[\widetilde g_i^2]\le4\cdot2^{-(2-\beta)w_i}
\le2^{-(2-2\beta)w_i}.
\]
The fact that $\operatorname{Var}[\widehat h_i]=\operatorname{Var}[\widehat g_i]=\mathbb E[\widehat g_i^2]$ finishes the proof.
\end{proof}
Now,
\begin{align*}
\mathbb E[\check S_{\ell,q}^2]
&=\sum_{\substack{I\subseteq B_q\\|I|=\ell}}\prod_{i\in I}\mathbb E[\widehat g_i^2]
\le\sum_{\substack{I\subseteq B_q\\|I|=\ell}}\prod_{i\in I}2^{-(2-2\beta)w_i}\\
&\le\sum_{\substack{I\subseteq B_q\\|I|=\ell}}2^{-(2-2\beta)c^q\ell}
\le\binom{m_q}{\ell}2^{-(2-2\beta)c^q\ell}\\
&\le\frac{3^{c^q\ell}e^\ell2^{-(2-2\beta)c^q\ell}}{\ell!}
\le\frac{2^{-c^q\ell/4}}{\ell!}.
\end{align*}
Plugging it in our expression for $R_{\vec k}$, we get
\begin{equation}
\mathbb E[R_{\vec k}(\vec g)]
=\prod_{q\in[Q]}\mathbb E[\check S_{\vec k[q],q}^2\vec k[q]!]
\le\prod_{q\in[Q]}2^{-c^q\vec k[q]/4}=2^{-\|\vec k\|_{(c)}/4}.\tag{13}
\end{equation}
Finally, let us bound $L_1(R_{\vec k}(\vec g))$. In Equation(12) we established the fact that $L_1(S_{\vec k}(\vec g))\le12^{\|\vec k\|_{(c)}}\le12^A$. Thus,
\begin{align*}
L_1(R_{\vec k}(\vec g))
&\le12^{2A}\prod_{q\in[Q]}\vec k[q]!
\le12^{2A}e^{\sum_{q\in[Q]}\vec k[q]\ln\vec k[q]}\\
&\le12^{2A}e^{(\ln A)\sum_{q\in[Q]}\vec k[q]}.
\end{align*}
As $\|\vec k\|_{(c)}=\sum_{q\in[Q]}c^q\vec k[q]\le A$ and $c^q\ge C\log\log n$ (see Claim18), $\sum_{q\in[Q]}\vec k[q]\le A/(C\log\log n)$ and we get
\[
L_1(R_{\vec k}(\vec g))
\le12^{2A}e^{\ln A\cdot A/(C\log\log n)}
\le12^{2A}2^{(\bar C/C)\log n}\le n^{10\bar C}.
\]
Note that $\delta\le\frac1{32}n^{-10\bar C}2^{-A/4}$. Thus, by Claim5,
\[
\mathbb E[R_{\vec k}(\vec g(Y))]
\le2^{-\|\vec k\|_{(c)}/4}+\delta\cdot n^{10\bar C}
\le2^{-\|\vec k\|_{(c)}/8},
\]
and we are done with bounding $\Pr[I^g(Y)=0]$. For the bound on $\mathbb E[S_{\vec k}^2(\vec g(Y))]$, simply observe that $\mathbb E[S_{\vec k}^2(\vec g(Y))]<\mathbb E[R_{\vec k}(\vec g(Y))]$.
\end{proof}

\section*{5 Full PRG via Iterated Restrictions}
So far, we have shown how to pseudorandomly assign values to a constant fraction of the inputs of any read-once PARITY $\circ$ AND formula using $O(\log n)$ truly random bits, preserving the expectation of the formula to within near-optimal error. In this section, to complete the proof of Theorem1, we show how to pseudorandomly assign values to all the inputs, i.e., we give a genuine PRG.
For convenience, we make the following definitions.
\begin{authorclaim}{Definition26}
Let $w>0$. A $w$-proper formula is a read-once PARITY $\circ$ AND formula of width at most $w$ and length most $2^{8w}$. We say that such a formula is short if its length is at most $2^{4w}$; otherwise, we say that the formula is long.
\end{authorclaim}
Our main goal is to fool $(C\log n)$-proper formulas, but along the way, we will obtain a PRG for $w$-proper formulas with seed length $O(w)$ and error $\exp(-\widetilde\Omega(w))$, even for $w$ substantially smaller than $\log n$.
\subsection*{5.1 Restrictions for Proper Formulas}
Recall that Lemma9 provides a pseudorandom restriction that uses only $O(\log n)$ truly random bits. We now generalize this fact in two respects. First, in the case of $w$-proper formulas ($\log\log n\le w\le C\log n$), we improve the seed length to $O(w)$. Second, in the case of short $w$-proper formulas, we argue that the restriction simplifies the formula, in the sense that it transforms it into a $(w/2)$-proper formula.
\begin{authorclaim}{Lemma27}
For every $w,n\in\mathbb N$ with $w\le C\log n$, there is a distribution $X$ over $\{0,1,?\}^n$ with the following properties.
\begin{enumerate}
\item (Seed length) There is an explicit algorithm to sample from $X$ using just $O(w+\log\log n)$ truly random bits.
\item (Expectation preservation) If $f$ is a $w$-proper formula, then $X$ preserves the expectation of $f$ with error $\exp(-\Omega(w/\log w))$.
\item (Simplification) If $f$ is a short $w$-proper formula, then $\Pr[f|_X\text{ is a }(w/2)\text{-proper formula}]\ge1-2^{-w}$.
\end{enumerate}
\end{authorclaim}
\begin{proof}
Let $n'=2^{8w}\cdot w$. Let $Y$ be an $n'$-wise $\delta^3$-biased distribution where $\delta=(n')^{-12\bar C}$, and let $Z$ be $\gamma$-almost $k$-wise independent with marginals $1-2^{-C}$, where $k=6\log n'$ and $\gamma=(n')^{-9}$. Our restriction is $X=\operatorname{Res}(Y,Z)^{\circ2^{C+4}}$.
By standard constructions [25,3] and Claim7, $X$ can be explicitly sampled using $O(w+\log\log n)$ truly random bits.
Now, to prove expectation preservation, let $f$ be a $w$-proper formula. By $w$-properness, there is some set of indices $I\subseteq[n]$, $|I|\le n'$, such that $f(x)$ only depends on $x|_I$. Let $g:\{0,1\}^{|I|}\to\{0,1\}$ be the $w$-proper formula such that $f(x)=g(x|_I)$. Since $Y|_I$ is $\delta^3$-biased and $Z|_I$ is $\gamma$-almost $k$-wise independent with marginals $1-2^{-C}$, Lemma9 implies that $\operatorname{Res}(Y|_I,Z|_I)$ preserves the expectation of $g$ with error $\exp(-\Omega(\log n'/\log\log n'))$, which is $\exp(-\Omega(w/\log w))$. It follows that $\operatorname{Res}(Y,Z)$ preserves the expectation of $f$ with the same error. The error of $X$ is only larger by a constant factor $2^{C+4}$, because any restriction of a $w$-proper formula is trivially another $w$-proper formula.
Finally, to prove simplification, let $f$ be a short $w$-proper formula, and let $f_i$ be a term. Since $k>w/2$, by Claim8, the probability that more than $w/2$ variables from $f_i$ are assigned ? by $X$ is bounded by
\begin{align*}
\binom w{w/2}\cdot((1-2^{-C})^{2^{C+4}\cdot w/2}+2^{C+4}\gamma)
&\le2^w\cdot(e^{-2^{-C}\cdot2^{C+3}w}+2^{C+4}\cdot(n')^{-9})\\
&<2^{-5w}.
\end{align*}
The number of terms in $f$ is at most $2^{4w}$, so by the union bound, except with probability $2^{-w}$, $f|_X$ has maximum width at most $w/2$. Furthermore, restricting cannot increase the number of terms, so the number of terms is still bounded by $2^{4w}=2^{8(w/2)}$. Therefore, in this case, $f|_X$ is $(w/2)$-proper.
\end{proof}
\subsection*{5.2 Full PRGs for Long Proper Formulas [24]}
The simplification clause of Lemma27 only applies if $f$ is short. If $f$ is long, we will therefore need a different approach. We will take a similar approach as Meka, Reingold, and Tal [24]. A full PRG for long $w$-proper formulas follows readily from their work.
\begin{authorclaim}{Lemma28 ([24])}
For every $w,n\in\mathbb N$, there is an explicit $(2^{-w})$-PRG for long $w$-proper formulas with seed length $O(w+\log\log n)$.
\end{authorclaim}
\begin{proof}[Proof sketch]
In short, the PRG is one of the PRGs by Meka et al. [24, full version, Lemma6.2], except we replace every $\delta$-biased distribution with a $(\cdot)$-wise $\delta$-biased distribution to optimize the seed length.
In more detail, let $n'=2^{8w}\cdot w$. Sample $v\in\{0,1\}^{wn}$ from an $(n'w)$-wise $(c_{\mathrm{MRT}}^{-w})$-biased distribution, where $c_{\mathrm{MRT}}$ is a suitable constant. Think of $v$ as $n$ blocks of $w$ bits. Define a set $I\subseteq[n]$ as follows: include $i$ in $I$ if and only if the $i$-th block of $v$ is $1^w$.
Sample $x^{(0)},x^{(1)},\ldots,x^{(16)}\in\{0,1\}^n$ independently from an $(n')$-wise $(c_{\mathrm{MRT}}^{-w})$-biased distribution. The PRG outputs the string $x$ defined by $x_i=x_i^{(0)}$ if $i\notin I$, and $x_i=\bigoplus_{j=1}^{16}x_i^{(j)}$ if $i\in I$.
By standard constructions [25,3], the seed length of this PRG is $O(\log n'+w+\log\log n)=O(w+\log\log n)$.
As for correctness, let $f$ be a long $w$-proper formula. Let $J\subseteq[n]$ be the set of indices of variables that $f$ reads, so there is some long $w$-proper formula $g$ on $|J|$ input bits such that $f(x)=g(x|_J)$. Let $X$ be the distribution output by the PRG. Since $|J|\le n'$, the distribution $X|_J$ is exactly the pseudorandom distribution designed by Meka et al. [24, full version, Lemma6.2]. Furthermore, since $f$ is long, $|J|>2^{4w}$. It follows that $g$ is in the class of functions fooled by Meka et al.'s pseudorandom distribution: $g$ is an XOR of $m$ non-constant Boolean functions on disjoint variables, where each function is on at most $w$ variables, with $16^w<m\le16^{2w}$ and $\log\log(|J|/2^w)\ll w\le\log|J|$. Therefore, $X|_J$ fools $g$ with error $2^{-w}$, and hence $X$ fools $f$ with error $2^{-w}$.
\end{proof}
\subsection*{5.3 Full PRGs for Width-$O(\log n)$ Formulas}
For short proper $w$-formulas, to get a full PRG, we will iterate the restriction of Lemma27 several times, assigning values to more and more variables. Eventually, we'll stop this recursive process and use a different PRG. Specifically, for the "base case," we'll use a PRG by Lee [22] with minor modifications:
\begin{authorclaim}{Lemma29 ([22])}
For every $w,n\in\mathbb N$ and every $\epsilon>0$, there is an explicit $\epsilon$-PRG for $w$-proper formulas with seed length
\[
O((w+\log(1/\epsilon))\cdot(\log w+\log\log(1/\epsilon))^2)
+\operatorname{poly}(\log\log(n/\epsilon)).
\]
\end{authorclaim}
\begin{proof}[Proof sketch]
In short, the PRG is one of the PRGs by Lee [22, Theorem6], except we replace every $\delta$-biased distribution with a $(\cdot)$-wise $\delta$-biased distribution to optimize the seed length, just like the proofs of Lemma27 and Lemma28.
To give a little more detail, let $n'=2^{8w}\cdot w$; a $w$-proper formula only reads $n'$ variables. Lee's PRG [22, Theorem6] is designed to fool arbitrary-order combinatorial checkerboards, i.e., parities of functions on disjoint variable sets of size at most $w$. This class includes $w$-proper formulas as a special case. Lee's original PRG has seed length $O((w+\log(n/\epsilon))\cdot(\log w+\log\log(n/\epsilon))^2)$.
After making suitable replacements, one can show that the seed length is reduced to
\[
O((w+\log(n'/\epsilon))\cdot(\log w+\log\log(n'/\epsilon))^2)
+\operatorname{poly}(\log\log(n/\epsilon)).
\]
(We omit the full proof, since it repeats much of Lee's analysis [22].) Plugging in the value of $n'$, we get the claimed seed length.
\end{proof}
We now give our full PRG for general formulas of width at most $C\log n$. The PRG follows a similar approach to one of the PRGs by Meka et al. [24, full version, Algorithm3]: iteratively apply the restriction of Lemma27, but at each step, XOR with the PRG of Lemma28 in case the formula is long.
\begin{authorclaim}{Lemma30}
For every $n\in\mathbb N$, there is an explicit PRG for read-once PARITY $\circ$ AND formulas of width at most $C\log n$ with seed length $O(\log n)$ and error $2^{-\Omega(\log n/(\log\log n)^3)}$.
\end{authorclaim}
\begin{proof}
Define $w_0=\log n/(\log\log n)^2$. We recursively define a PRG $G_w$ for $w$-proper formulas, $w_0\le w\le C\log n$, as follows.
\begin{itemize}
\item (Base case) If $w\le2w_0$, then $G_w$ is the $(2^{-w_0})$-PRG of Lemma29 based on Lee's work [22].
\item (Recursive case) If $w>2w_0$, sample $X\in\{0,1,?\}^n$ from the distribution guaranteed by Lemma27 based on the work in Section4. Sample $Y\in\{0,1\}^n$ using the PRG of Lemma28 based on Meka et al.'s work [24]. Recursively sample $G_{\lceil w/2\rceil}$, and set $G_w=Y\oplus(X\circ G_{\lceil w/2\rceil})$.
\end{itemize}
For the analysis, observe first that in the base case $w\le2w_0$, $G_w$ fools $w$-proper formulas with error $2^{-w_0}$. Now, for the inductive step, consider some $w>2w_0$. Assume $G_{\lceil w/2\rceil}$ fools $\lceil w/2\rceil$-proper formulas with error $\epsilon_{\lceil w/2\rceil}$; we will show that $G_w$ fools $w$-proper formulas with error $\epsilon_w$, where $\epsilon_w=\epsilon_{\lceil w/2\rceil}+2^{-\Omega(w/\log w)}$.
Let $f$ be a $w$-proper formula, and for brevity, let $G=G_{\lceil w/2\rceil}$. For the first case, suppose $f$ is long. Any shift of $f$ is also a long $w$-proper formula, so
\begin{align*}
|\mathbb E[f(G_w)]-\mathbb E[f]|
&=\left|\mathbb E_{X,G}[\mathbb E_Y[f(Y\oplus(X\circ G))]]-\mathbb E[f]\right|\\
&\le\mathbb E_{X,G}[|\mathbb E_Y[f(Y\oplus(X\circ G))]-\mathbb E[f]|]\\
&=\mathbb E_{X,G}[|\mathbb E_Y[f(Y\oplus(X\circ G))]-\mathbb E_U[f(U\oplus(X\circ G))]|]\\
&\le2^{-w}.
\end{align*}
For the second case, suppose $f$ is short. For each $y\in\{0,1\}^n$, define $f_y(x)=f(y\oplus x)$, another short $w$-proper formula. Fix $y\sim Y$, and let $\mathcal E$ be the event that $f_y|_X$ is $(w/2)$-proper, so whether $\mathcal E$ occurs depends only on $X$. Then
\begin{align*}
|\mathbb E[(f_y|_X)(G)]-\mathbb E[f]|
&\le|\mathbb E_X[\mathbb E_G[(f_y|_X)(G)]\mid\mathcal E]-\mathbb E[f]|+\Pr[\neg\mathcal E]\\
&\le|\mathbb E_X[\mathbb E_U[(f_y|_X)(U)]\mid\mathcal E]-\mathbb E[f]|
+\epsilon_{\lceil w/2\rceil}+\Pr[\neg\mathcal E]\quad(\text{Induction})\\
&\le|\mathbb E_X[\mathbb E_U[(f_y|_X)(U)]]-\mathbb E[f]|
+\epsilon_{\lceil w/2\rceil}+2\Pr[\neg\mathcal E]\\
&\le2^{-\Omega(w/\log w)}+\epsilon_{\lceil w/2\rceil}+2\Pr[\neg\mathcal E]\quad(\text{Item2 of Lemma27})\\
&\le2^{-\Omega(w/\log w)}+\epsilon_{\lceil w/2\rceil}+2\cdot2^{-w}\quad(\text{Item3 of Lemma27}).
\end{align*}
Let $\epsilon_w$ be the final right-hand side, so indeed $\epsilon_w=\epsilon_{\lceil w/2\rceil}+\exp(-\Omega(w/\log w))$. Then
\[
|\mathbb E[f(G_w)]-\mathbb E[f]|
\le\mathbb E_Y[|\mathbb E_{X,G}[(f_Y|_X)(G)]-\mathbb E[f]|]\le\epsilon_w.
\]
Now, let us add up all these errors. Since $w\ge w_0$ always holds, we have $\epsilon_w\le\epsilon_{\lceil w/2\rceil}+\exp(-\Omega(w_0/\log w_0))$. Starting at $w=C\log n$, we only need to halve $w$ a total of $O(\log\log\log n)$ times to reach the base case $w\le2w_0$. Therefore, the total error of $G_{C\log n}$ is bounded by
\[
2^{-w_0}+2^{-\Omega(w_0/\log w_0)}\cdot O(\log\log\log n)
=2^{-\Omega(\log n/(\log\log n)^3)}.
\]
Finally, let us bound the seed length of $G_w$. In the base case $w\le2w_0$, by our choice of $w_0$, the seed length $s_w$ of $G_w$ is bounded by some value $s_{\mathrm{base}}\le O(\log n)$. In the recursive case $w>2w_0$, the seed length $s_w$ of $G_w$ is bounded by
$s_w=s_{\lceil w/2\rceil}+O(w+\log\log n)=s_{\lceil w/2\rceil}+O(w)$.
The point is that this is essentially a geometric series. More precisely, let $c_{\mathrm{seed}}$ be a constant such that $s_w\le s_{\lceil w/2\rceil}+c_{\mathrm{seed}}\cdot w$ for all $w>2w_0$. Then by induction, for all $w\ge w_0$, we have $s_w\le s_{\mathrm{base}}+3c_{\mathrm{seed}}w$, because
\begin{align*}
s_w&\le s_{\lceil w/2\rceil}+c_{\mathrm{seed}}w\\
&\le s_{\mathrm{base}}+3c_{\mathrm{seed}}\lceil w/2\rceil+c_{\mathrm{seed}}w\quad(\text{Induction})\\
&<s_{\mathrm{base}}+3c_{\mathrm{seed}}w.
\end{align*}
Therefore, we can take the desired PRG to be $G_{C\log n}$, because $s_{C\log n}\le O(\log n)$, and any read-once PARITY $\circ$ AND formula of width at most $C\log n$ is $(C\log n)$-proper.
\end{proof}
\subsection*{5.4 Arbitrary-Error PRGs for Width-$O(\log(n/\epsilon))$ Formulas}
At this point, the main work of proving Theorem1 is complete. We just need to address three minor issues: small $\epsilon$, large width, and formulas not of the form PARITY $\circ$ AND. We begin by addressing the case of small $\epsilon$. Recall that we wish to achieve seed length $O(\log n)+\widetilde O(\log(1/\epsilon))$ for an arbitrary error $\epsilon$. This follows readily by combining the PRG of Lemma30 with Lee's PRG (Lemma29).
\begin{authorclaim}{Lemma31}
For any $n\in\mathbb N$, $\epsilon>0$, there is an explict $\epsilon$-PRG for read-once PARITY $\circ$ AND formulas of width at most $(C/2)\log(n/\epsilon)$ with seed length $O(\log n+\log(1/\epsilon)\cdot(\log\log(1/\epsilon))^5)$.
\end{authorclaim}
\begin{proof}
Let $\epsilon_0$ be the error parameter in Lemma30, so $\epsilon_0=\exp(-\Omega(\log n/(\log\log n)^3))$. If $\epsilon\ge\epsilon_0$, the PRG of Lemma30 works, because $(C/2)\log(n/\epsilon)<C\log n$. If $\epsilon<\epsilon_0$, use Lee's PRG [22], i.e., the $\epsilon$-PRG of Lemma29 for $((C/2)\log(n/\epsilon))$-proper formulas, which has seed length
\[
O(\log(n/\epsilon)\cdot(\log\log(n/\epsilon))^2)
\le O(\log(1/\epsilon)\cdot(\log\log(1/\epsilon))^5).
\]
\end{proof}
(In the proof of Lemma31, we could just as well have used Lee's original PRG [22, Theorem6] instead of the slightly modified version given by Lemma29.)
\subsection*{5.5 PRGs for Any Width}
In this section, we eliminate the assumption that the maximum width is bounded.
\begin{authorclaim}{Lemma32}
For all $n\in\mathbb N$ and $\epsilon>0$, there is an explicit $\epsilon$-PRG for read-once PARITY $\circ$ AND formulas on $n$ input bits with seed length $O(\log n+\log(1/\epsilon)\cdot(\log\log(1/\epsilon))^5)$.
\end{authorclaim}
\begin{proof}
Sample $G$ from the $(\epsilon/3)$-PRG for formulas of width $(C/2)\log(3n/\epsilon)$ guaranteed by Lemma31. Sample $Y$ from an $(\epsilon/(6n))$-biased distribution. Our final PRG outputs $H\mathrel{\mathop:}=G\oplus Y$.
To prove that this works, let $f$ be a read-once PARITY $\circ$ AND formula. Write $f=f'\oplus f''$, where every term in $f'$ has width at most $(C/2)\log(3n/\epsilon)$ and every term in $f''$ has width greater than $(C/2)\log(3n/\epsilon)$.
Since any shift of a width-$w$ read-once PARITY $\circ$ AND formula is another width-$w$ read-once PARITY $\circ$ AND formula, $H$ fools $f'$ with error $\epsilon/3$. Meanwhile, since each term $f''_i$ of $f''$ is a conjunction of more than $(C/2)\log(3n/\epsilon)$ literals, $\mathbb E[f''_i]\le(\epsilon/(3n))^{C/2}<\epsilon/(6n)$.
Furthermore, the $L_1$ norm of any conjunction of literals is1, and $H$ is $(\epsilon/(6n))$-biased, so by Claim5, $\mathbb E[f''_i(H)]<\epsilon/(3n)$. Therefore, by the union bound, for either distribution $X\in\{H,U\}$, $\mathbb E[f''(X)]<\epsilon/3$.
This allows us to bound the error of the final PRG as follows:
\begin{align*}
|\mathbb E[f(H)]-\mathbb E[f]|
&\le|\mathbb E[f(H)]-\mathbb E[f'(H)]|+|\mathbb E[f'(H)]-\mathbb E[f']|+|\mathbb E[f']-\mathbb E[f]|\\
&\le\mathbb E[|f(H)-f'(H)|]+|\mathbb E[f'(H)]-\mathbb E[f']|+\mathbb E[|f'-f|]\\
&=\mathbb E[f''(H)]+|\mathbb E[f'(H)]-\mathbb E[f']|+\mathbb E[f'']\\
&<\epsilon/3+\epsilon/3+\epsilon/3=\epsilon.
\end{align*}
\end{proof}
\subsection*{5.6 Proof of Theorem1}
In this section, we finally complete the proof of Theorem1 by showing that fooling read-once PARITY $\circ$ AND formulas is sufficient for fooling read-once depth-2 $\mathrm{AC}^0[\oplus]$:
\begin{authorclaim}{Lemma33}
Let $X$ be a distribution over $\{0,1\}^n$, and let $\epsilon>0$. If $X$ fools all read-once PARITY $\circ$ AND formulas with error $\epsilon$, then $X$ fools all read-once depth-2 $\mathrm{AC}^0[\oplus]$ formulas with error $2\epsilon$.
\end{authorclaim}
\begin{proof}
Let $f$ be a read-once depth-2 $\mathrm{AC}^0[\oplus]$ formula.
For the first case, suppose the output gate of $f$ is $\oplus$. By merging the output gate with any $\oplus$ children and introducing trivial $\wedge$ gates with fan-in1 as necessary, we see that without loss of generality, every child of the output gate is either $\wedge$ or $\vee$. By de Morgan's laws, it follows that either $f$ or $\neg f$ can be computed by a read-once PARITY $\circ$ AND formula. Either way, this implies that $X$ $\epsilon$-fools $f$.
For the second case, suppose the output gate of $f$ is $\wedge$, say $f=\bigwedge_{i=1}^mf_i$. Using the Fourier expansion of the $m$-input AND function, we get
\[
f=\sum_{I\subseteq[m]}\frac{(-1)^{|I|}}{2^m}\prod_{i\in I}(-1)^{f_i}
=\sum_{I\subseteq[m]}\frac{(-1)^{|I|}}{2^m}\cdot\left(1-2\cdot\bigoplus_{i\in I}f_i\right).
\]
By our analysis for the first case, $X$ fools $\bigoplus_{i\in I}f_i$ with error $\epsilon$. Therefore, by the triangle inequality,
\begin{align*}
|\mathbb E[f(X)]-\mathbb E[f]|
&\le\left|\sum_{I\subseteq[m]}\frac{(-1)^{|I|}\cdot(-2)}{2^m}
\left(\mathbb E[(\bigoplus_{i\in I}f_i)(X)]-\mathbb E[\bigoplus_{i\in I}f_i]\right)\right|\\
&\le\sum_{I\subseteq[m]}\frac2{2^m}\cdot\epsilon=2\epsilon.
\end{align*}
For the final case, suppose the output gate of $f$ is $\vee$. By de Morgan's laws, $\neg f$ can be computed by a read-once depth-2 $\mathrm{AC}^0[\oplus]$ formula with output gate $\wedge$. By our analysis for the second case, $X$ fools $\neg f$ with error $2\epsilon$, hence $X$ fools $f$ with the same error.
\end{proof}

\begin{thebibliography}{33}
\bibitem{r1} M. Ajtai, J. Komlos, and E. Szemeredi. Deterministic simulation in LOGSPACE. In Proceedings of the Nineteenth Annual ACM Symposium on Theory of Computing (STOC), pages 132–140, New York, NY, USA, 1987. ACM. doi:10.1145/28395.28410.
\bibitem{r2} Miklos Ajtai and Avi Wigderson. Deterministic simulation of probabilistic constant depth circuits. Advances in Computing Research, 5(199-222):1, 1989.
\bibitem{r3} Noga Alon, Oded Goldreich, Johan Håstad, and René Peralta. Simple constructions of almost k-wise independent random variables. Random Structures \& Algorithms, 3(3):289–304, 1992.
\bibitem{r4} Louay Bazzi and Nagi Nahas. Small-bias is not enough to hit read-once CNF. Theory of Computing Systems, 60(2):324–345, February 2017. doi:10.1007/s00224-016-9680-6.
\bibitem{r5} Mark Braverman, Gil Cohen, and Sumegha Garg. Pseudorandom pseudo-distributions with near-optimal error for read-once branching programs. SIAM Journal on Computing, 0(0):STOC18–242–STOC18–299, 2020. doi:10.1137/18M1197734.
\bibitem{r6} Mark Braverman, Anup Rao, Ran Raz, and Amir Yehudayoff. Pseudorandom generators for regular branching programs. SIAM Journal on Computing, 43(3):973–986, 2014.
\bibitem{r7} J. Brody and E. Verbin. The coin problem and pseudorandomness for branching programs. In 2010 IEEE 51st Annual Symposium on Foundations of Computer Science (FOCS), pages 30–39, October 2010. doi:10.1109/FOCS.2010.10.
\bibitem{r8} Suresh Chari, Pankaj Rohatgi, and Aravind Srinivasan. Improved algorithms via approximations of probability distributions. Journal of Computer and System Sciences, 61(1):81–107, 2000. doi:10.1006/jcss.1999.1695.
\bibitem{r9} Eshan Chattopadhyay, Pooya Hatami, Omer Reingold, and Avishay Tal. Improved pseudorandomness for unordered branching programs through local monotonicity. In Proceedings of the 50th Annual ACM Symposium on Theory of Computing (STOC), pages 363–375, New York, NY, USA, 2018. ACM. doi:10.1145/3188745.3188800.
\bibitem{r10} Sitan Chen, Thomas Steinke, and Salil Vadhan. Pseudorandomness for read-once, constantdepth circuits. arXiv preprint, 2015. arXiv:1504.04675.
\bibitem{r11} Anindya De. Pseudorandomness for permutation and regular branching programs. In Proceedings of the 26th Annual IEEE 26th Annual Conference on Computational Complexity (CCC), pages 221–231. IEEE, 2011.
\bibitem{r12} Anindya De, Omid Etesami, Luca Trevisan, and Madhur Tulsiani. Improved pseudorandom generators for depth 2 circuits. In Approximation, randomization, and combinatorial optimization, volume 6302 of Lecture Notes in Computer Science, pages 504–517. Springer, Berlin, 2010. doi:10.1007/978-3-642-15369-3\_38.
\bibitem{r13} Dean Doron, Pooya Hatami, and William M Hoza. Near-optimal pseudorandom generators for constant-depth read-once formulas. In 34th Computational Complexity Conference (CCC), 2019.
\bibitem{r14} Michael A. Forbes and Zander Kelley. Pseudorandom generators for read-once branching programs, in any order. In Proceedings of the 59th Annual IEEE Symposium on Foundations of Computer Science (FOCS). IEEE, 2018.
\bibitem{r15} Anat Ganor and Ran Raz. Space pseudorandom generators by communication complexity lower bounds. In LIPIcs-Leibniz International Proceedings in Informatics, volume 28. Schloss Dagstuhl-Leibniz-Zentrum fuer Informatik, 2014.
\bibitem{r16} Parikshit Gopalan, Raghu Meka, Omer Reingold, Luca Trevisan, and Salil Vadhan. Better pseudorandom generators from milder pseudorandom restrictions. In Proceedings of the 53rd Annual IEEE Symposium on Foundations of Computer Science (FOCS), pages 120–129. IEEE, 2012.
\bibitem{r17} Parikshit Gopalan and Amir Yehudayoff. Inequalities and tail bounds for elementary symmetric polynomial with applications. arXiv preprint, 2014. arXiv:1402.3543.
\bibitem{r18} Elad Haramaty, Chin Ho Lee, and Emanuele Viola. Bounded independence plus noise fools products. SIAM Journal on Computing, 47(2):493–523, 2018. doi:10.1137/17M1129088.
\bibitem{r19} William M. Hoza and David Zuckerman. Simple optimal hitting sets for small-success RL. In Proceedings of the 59th Annual IEEE Symposium on Foundations of Computer Science (FOCS). IEEE, 2018.
\bibitem{r20} Russell Impagliazzo, Noam Nisan, and Avi Wigderson. Pseudorandomness for network algorithms. In Proceedings of the Twenty-Sixth Annual ACM Symposium on Theory of Computing (STOC), pages 356–364. ACM, 1994.
\bibitem{r21} Michal Koucký, Prajakta Nimbhorkar, and Pavel Pudlák. Pseudorandom generators for group products. In Proceedings of the 43rd Annual ACM Symposium on Theory of Computing (STOC), pages 263–272. ACM, New York, 2011. doi:10.1145/1993636.1993672.
\bibitem{r22} Chin Ho Lee. Fourier bounds and pseudorandom generators for product tests. In Amir Shpilka, editor, 34th Computational Complexity Conference (CCC), volume 137 of Leibniz International Proceedings in Informatics (LIPIcs), pages 7:1–7:25, Dagstuhl, Germany, 2019. Schloss Dagstuhl–Leibniz-Zentrum fuer Informatik. doi:10.4230/LIPIcs.CCC.2019.7.
\bibitem{r23} Chin Ho Lee and Emanuele Viola. More on bounded independence plus noise: Pseudorandom generators for read-once polynomials. In Electronic Colloquium on Computational Complexity (ECCC), volume 24, page 167, 2017.
\bibitem{r24} Raghu Meka, Omer Reingold, and Avishay Tal. Pseudorandom generators for width-3 branching programs. In Proceedings of the 51st Annual ACM Symposium on Theory of Computing (STOC), pages 626–637. ACM, New York, 2019.
\bibitem{r25} Joseph Naor and Moni Naor. Small-bias probability spaces: Efficient constructions and applications. SIAM Journal on Computing, 22(4):838–856, 1993.
\bibitem{r26} Noam Nisan. Pseudorandom generators for space-bounded computation. Combinatorica, 12(4):449–461, 1992.
\bibitem{r27} Noam Nisan. $\mathrm{RL}\subseteq\mathrm{SC}$. computational complexity, 4(1):1–11, March 1994. doi:10.1007/ BF01205052.
\bibitem{r28} Noam Nisan and David Zuckerman. Randomness is linear in space. Journal of Computer and System Sciences, 52(1):43–52, 1996. doi:10.1006/jcss.1996.0004.
\bibitem{r29} Omer Reingold. Undirected connectivity in log-space. Journal of the ACM, 55(4):Art. 17, 24, 2008. doi:10.1145/1391289.1391291.
\bibitem{r30} Omer Reingold, Thomas Steinke, and Salil Vadhan. Pseudorandomness for regular branching programs via Fourier analysis. In Approximation, Randomization, and Combinatorial Optimization. Algorithms and Techniques, pages 655–670. Springer, 2013.
\bibitem{r31} Michael E. Saks and Shiyu Zhou. $\mathrm{BPHSPACE}(S)\subseteq\mathrm{DSPACE}(S^{3/2})$. Journal of Computer and System Sciences, 58(2):376–403, 1999.
\bibitem{r32} Thomas Steinke. Pseudorandomness for permutation branching programs without the group theory. In Electronic Colloquium on Computational Complexity (ECCC), 2012.
\bibitem{r33} Thomas Steinke, Salil Vadhan, and Andrew Wan. Pseudorandomness and Fourier-growth bounds for width-3 branching programs. Theory of Computing, 13(12):1–50, 2017. doi: 10.4086/toc.2017.v013a012.
\end{thebibliography}
\end{document}
